\chapter{Homotopy invariance of the index and Toeplitz operators}\label{ch:index}
\chaptermark{Index and Toeplitz operators}
\providecommand{\Ginf}{G^{-\infty}}
\providecommand{\Gloop}{G^{-\infty}_{(1)}}
\providecommand{\PsiT}[1]{\Psi^{#1}_{\mathcal T}}
\providecommand{\Hardy}{\mathcal{C}^\infty_{+}}
\providecommand{\Hardym}{\mathcal{C}^\infty_{-}}
\providecommand{\wind}{\operatorname{wind}}
\providecommand{\Iso}{\operatorname{Iso}}

An elliptic operator on a compact manifold is Fredholm: its null space and its
cokernel are finite dimensional. Neither dimension is stable. A small
perturbation can create or destroy solutions of $Au=0$. Their difference, the
index, does not move. This stability is the reason the index can be computed
from topological data, and this chapter explains where it comes from. It also
builds the first tools for the computation: K-theory, Chern forms and the
simplest index formula, for Toeplitz operators on the circle.
\Cref{index:fig:roadmap} shows how the sections fit together.

This chapter opens the part of the book devoted to index theory. Let $X$ be a
compact manifold without boundary and $E$, $F$ complex vector bundles over $X$.
By \cref{ch:manifolds} an elliptic operator $A\in\Psido{m}(X;E,F)$ has a
parametrix and is Fredholm, so its index
\begin{equation*}
  \ind(A)=\dim\Nul(A)-\dim\bigl(\CI(X;F)/A\,\CI(X;E)\bigr)
\end{equation*}
is defined. The first two sections show that $\ind(A)$ can be written as a
difference of traces of smoothing operators, and deduce from this that the
index is constant along smooth families of elliptic operators. It therefore
depends only on the homotopy class of the principal symbol. The index theorem
computes this integer from topological data. To set that up we take
\emph{odd K-theory} to be the homotopy classes of maps into the group
$\Ginf(Y;E)$ of invertible smoothing perturbations of the identity. We show that
this is an abelian group and introduce its even counterpart, which is
classified by a loop group $\Gloop$ (\cref{index:sec:K}). The bi-invariant
\emph{Chern forms} $\Tr((A^{-1}dA)^{2k-1})$ on $\Ginf$ give de~Rham classes on
these K-groups. The simplest of them, $\Tr(A^{-1}dA)$, has integral periods,
and this gives a direct construction of the Fredholm determinant
(\cref{index:sec:chern}). The last three sections turn to the Toeplitz algebra
on the circle. It is the model for the ``symbol sequence'' that leads to Bott
periodicity in \cref{ch:bott}. We compute the index of Toeplitz operators. We
show that families of elliptic operators with vanishing ``families index'' can
be perturbed to be invertible, and that the matrix groups can be deformed away
inside the invertible Toeplitz operators.

\begin{figure}[ht]
\centering
\begin{tikzpicture}[
  box/.style={draw,rounded corners,align=center,font=\small,
    inner sep=4pt,text width=3.4cm},
  >=Stealth]
\node[box] (tr) at (0,0) {index as a trace\\ \cref{index:sec:trace}};
\node[box] (ho) at (4.6,0) {homotopy invariance\\ \cref{index:sec:homotopy}};
\node[box] (sy) at (9.2,0) {index depends only on\\ the class of $\sigma_m(A)$};
\node[box] (K)  at (9.2,-2) {K-theory via $\Ginf$\\ and loops, \cref{index:sec:K}};
\node[box] (ch) at (4.6,-2) {Chern forms,\\ determinant, \cref{index:sec:chern}};
\node[box] (to) at (0,-2) {Toeplitz index\\ $=-\wind$, \cref{index:sec:toeplitz}};
\node[box] (fa) at (0,-4) {perturbing families,\\ \cref{index:sec:families}};
\node[box] (de) at (4.6,-4) {deforming $\GL(N,\C)$\\ away, \cref{index:sec:deform}};
\node[box,dashed] (bo) at (9.2,-4) {Bott periodicity\\ \cref{ch:bott}};
\draw[->] (tr)--(ho);
\draw[->] (ho)--(sy);
\draw[->] (sy)--node[right,font=\footnotesize]{organize classes}(K);
\draw[->] (K)--(ch);
\draw[->] (ch)--node[above,font=\footnotesize]{$u_1$}(to);
\draw[->] (to)--(fa);
\draw[->] (fa)--(de);
\draw[->] (de)--(bo);
\end{tikzpicture}
\caption{The logic of \cref{ch:index}. The trace formula gives homotopy
invariance; homotopy invariance leads to K-theory; the first Chern form $u_1$
computes the Toeplitz index as a winding number; the last two sections prepare the
proof of Bott periodicity in \cref{ch:bott}.}
\label{index:fig:roadmap}
\end{figure}

\section{The index as a trace}\label{index:sec:trace}

The index is the difference of two dimensions. Dimensions of finite rank
projections are traces (\cref{index:lem:proj-trace}), and traces of smoothing
operators depend continuously on the operator. So the first step is to write
the index as a difference of traces of smoothing operators built from a
parametrix.

\subsection{What we use from earlier chapters}
Throughout, $X$ is a compact manifold without boundary, of positive dimension,
and $E$, $F$ are complex vector bundles over it. Operators in
$\Psido{m}(X;E,F)$ are classical (one-step polyhomogeneous). As in
\cref{ch:conormal,ch:manifolds} their Schwartz kernels are conormal
distributions at the diagonal. The principal symbol $\sigma_m(A)$ is a
section of $\hom(\pi^*E,\pi^*F)$ over $T^*X\setminus0$, homogeneous of degree
$m$. After fixing a Riemannian metric we identify it with its restriction to
the cosphere bundle $S^*X$. We use the following facts.
\begin{enumerate}[label=(\roman*)]
\item The smoothing operators are
  $\Psido{-\infty}(X;E,F)=\CI(X^2;\Hom(E,F)\otimes\Omega_R)$, with
  $\Hom(E,F)=\pi_L^*F\otimes\pi_R^*E'$ and $\Omega_R$ the density bundle
  from the right factor. They carry the $\CI$ (Fr\'echet) topology. A
  smoothing operator maps $\CmI(X;E)$ continuously into $\CI(X;F)$. In
  particular it is bounded on $L^2$ for any choice of Hermitian metrics and of a
  smooth positive density, and the operator norm is a continuous seminorm on
  $\Psido{-\infty}$.
\item If $A\in\Psido{m}$ and $R\in\Psido{-\infty}$ then $AR$ and $RA$ are
  smoothing (when defined), and depend continuously on $R$. The formal
  adjoint and the transpose of $A$ lie in $\Psido{m}$, and
  $\sigma_m(A^*)=\sigma_m(A)^*$.
\item $A\in\Psido{m}(X;E,F)$ is bounded $\Hs{s}(X;E)\to\Hs{s-m}(X;F)$ for
  every $s$. If $A$ is elliptic, i.e.\ $\sigma_m(A)$ is invertible on $S^*X$,
  it has a parametrix $B_0\in\Psido{-m}(X;F,E)$:
  \begin{equation}\label{index:eq:parametrix}
    B_0A=\Id_E-R,\quad AB_0=\Id_F-R',\quad R\in\Psido{-\infty}(X;E),\
    R'\in\Psido{-\infty}(X;F).
  \end{equation}
\item Finite rank smoothing operators are dense in $\Psido{-\infty}(X;E,F)$.
  A smoothing operator of finite rank has kernel
  $\sum_{i=1}^Nf_i(x)e_i(y)$ with $f_i\in\CI(X;F)$ and
  $e_i\in\CI(X;E'\otimes\Omega)$. The approximation can be made uniform for
  families: if $a\in\CI(Z;\Psido{-\infty}(X;E,F))$ with $Z$ compact, the
  approximating kernels may be taken as $\sum_{i,j\le N}c_{ij}(z)f_i(x)e_j(y)$
  with fixed $f_i$, $e_j$.
\item The trace $\Tr(A)=\int_X\tr A(x,x)$ is a continuous linear functional on
  $\Psido{-\infty}(X;E)$. It equals the matrix trace on finite rank operators,
  and $\Tr(AB)=\Tr(BA)$ for $A\in\Psido{-\infty}(X;E,F)$,
  $B\in\Psido{-\infty}(X;F,E)$. The proof in \cref{ch:manifolds} for $E=F$
  applies verbatim.
\item If $A\in\Psido{-\infty}(X;E)$ and $\Id+A$ is invertible on $L^2(X;E)$
  (equivalently on $\CI(X;E)$), then $(\Id+A)^{-1}=\Id+B$ with
  $B\in\Psido{-\infty}(X;E)$. The group
  \begin{equation*}
    \Ginf(X;E)=\{\Id+A;\ A\in\Psido{-\infty}(X;E),\ \Id+A\text{ invertible}\}
  \end{equation*}
  is an open subset of $\Id+\Psido{-\infty}(X;E)$, a topological group, and
  connected.
\end{enumerate}
We identify $\Ginf(X;E)$ with the open subset
$\{A;\ \Id+A\text{ invertible}\}\subset\Psido{-\infty}(X;E)$ whenever we
speak of smooth maps into it. We also need two simple facts about finite rank
operators.

\begin{lemma}\label{index:lem:proj-trace}
  If $\pi\in\Psido{-\infty}(X;E)$ is a projection ($\pi^2=\pi$) of finite rank,
  then $\Tr(\pi)=\dim\Ran(\pi)$.
\end{lemma}
\begin{proof}
  Let $f_1,\dots,f_N\in\CI(X;E)$ be a basis of $\Ran(\pi)$. Then
  $\pi u=\sum_i\ell_i(u)f_i$, where the $\ell_i$ are continuous linear
  functionals given by smooth sections $e_i$ of $E'\otimes\Omega$. Since
  $\pi f_j=f_j$ we get $\ell_i(f_j)=\delta_{ij}$, so
  $\Tr(\pi)=\sum_i\ell_i(f_i)=N$.
\end{proof}

\begin{lemma}\label{index:lem:trace-commute}
  If $A\in\Psido{m}(X;E,F)$ and $C\in\Psido{-\infty}(X;F,E)$ then
  $\Tr(AC)=\Tr(CA)$, the traces being taken on $\Psido{-\infty}(X;F)$ and
  $\Psido{-\infty}(X;E)$ respectively.
\end{lemma}
\begin{proof}
  Both sides are continuous in $C$ by (ii) and (v). By (iv) it suffices to
  treat $C$ of rank one, $Cf=\ell(f)w$ with $w\in\CI(X;E)$ and
  $\ell(f)=\int_Xv\cdot f$, $v\in\CI(X;F'\otimes\Omega)$. Then $AC$ is the rank
  one operator $f\mapsto\ell(f)\,Aw$ and $CA$ is $g\mapsto\ell(Ag)\,w$. The
  functional $g\mapsto\ell(Ag)$ is given by the smooth section $A^tv$, where
  $A^t\in\Psido{m}$ is the transpose. Thus
  \begin{equation*}
    \Tr(AC)=\int_Xv\cdot Aw=\int_X(A^tv)\cdot w=\Tr(CA).\qedhere
  \end{equation*}
\end{proof}

\subsection{The generalized inverse}
The parametrix of an elliptic operator can be improved to an exact inverse up to
finite rank projections.

\begin{proposition}\label{index:prop:generalized-inverse}
  Let $A\in\Psido{m}(X;E,F)$ be elliptic. Choose Hermitian metrics on $E$, $F$
  and a smooth positive density on $X$, and let $A^*\in\Psido{m}(X;F,E)$ be
  the formal adjoint. Then $\Nul(A)\subset\CI(X;E)$ and
  $\Nul(A^*)\subset\CI(X;F)$ are finite dimensional and
  \begin{equation}\label{index:eq:range}
    \CI(X;F)=A\bigl(\CI(X;E)\bigr)\oplus\Nul(A^*),
  \end{equation}
  an $L^2$-orthogonal decomposition. Let $\pi$ and $\pi'$ be the orthogonal
  projections onto $\Nul(A)$ and $\Nul(A^*)$. There is a unique
  $B\in\Psido{-m}(X;F,E)$, the \emph{generalized inverse} of $A$, such that
  \begin{equation}\label{index:eq:gen-inverse}
    BA=\Id_E-\pi,\quad AB=\Id_F-\pi',\quad \pi B=0,\quad B\pi'=0.
  \end{equation}
  The projections $\pi,\pi'$ are self-adjoint finite rank smoothing
  operators, and
  \begin{equation*}
    \ind(A)=\dim\Nul(A)-\dim\Nul(A^*)=\Tr(\pi)-\Tr(\pi').
  \end{equation*}
\end{proposition}
\begin{proof}
  Let $B_0$, $R$, $R'$ be as in \eqref{index:eq:parametrix}. As a bounded map
  $A:\Hs{m}(X;E)\to L^2(X;F)$, $A$ satisfies $B_0A=\Id-R$ and $AB_0=\Id-R'$,
  where $R$ is compact on $\Hs m$ and $R'$ is compact on $L^2$. Indeed both
  map into $\CI$ continuously from $\Hs{-N}$ for every $N$, and
  $\Hs{N}\hookrightarrow\Hs{m}$ is compact for $N>m$ (Rellich). By
  \cref{ch:hilbert} $A:\Hs{m}\to L^2$ is Fredholm. In particular its range is
  closed. A distribution $u\in\Hs m$ with $Au=0$ satisfies $u=Ru\in\CI$, so
  $\Nul(A)$ is the same on $\Hs m$ and on $\CI$, and is finite
  dimensional. The orthogonal complement of $A(\Hs m)$ in $L^2$ consists of
  the $v\in L^2$ with $A^*v=0$ in the sense of distributions. Since
  $B_0^*A^*=\Id-(R')^*$, such $v$ satisfy $v=(R')^*v\in\CI$, so this
  complement is $\Nul(A^*)$, a finite dimensional subspace of $\CI(X;F)$. If
  $f\in\CI(X;F)$ is orthogonal to $\Nul(A^*)$ then $f=Au$ with $u\in\Hs m$,
  and $u=B_0f+Ru\in\CI$. This proves \eqref{index:eq:range}. The orthogonal
  projection onto a finite dimensional space $W\subset\CI$ with orthonormal
  basis $w_i$ has kernel $\sum_iw_i(x)\overline{w_i(y)}\nu(y)$, so $\pi$ and
  $\pi'$ are self-adjoint finite rank smoothing operators.

  The map $A:\Nul(A)^\perp\cap\Hs m\to A(\Hs m)$ is a continuous bijection
  between Hilbert spaces. Let $B:L^2(X;F)\to\Hs m(X;E)$ be its inverse on
  $A(\Hs m)$, extended by zero on $\Nul(A^*)$. By the open mapping theorem $B$
  is bounded, and \eqref{index:eq:gen-inverse} holds on $L^2$. These four
  identities determine $B$. Substituting
  $\Id=AB_0+R'$ on the right and $\Id=B_0A+R$ on the left,
  \begin{align*}
    B&=B(AB_0+R')=(\Id-\pi)B_0+BR'
     =(\Id-\pi)B_0+(B_0A+R)BR'\\
     &=(\Id-\pi)B_0+B_0(\Id-\pi')R'+RBR'.
  \end{align*}
  The first two terms lie in $\Psido{-m}(X;F,E)$. The kernel of $RBR'$ is
  smooth. The map $y\mapsto R'(\cdot,y)$ is smooth from $X$ into $\CI(X;F)$,
  hence into $L^2$. Then $B$ is bounded on $L^2$, and $R:L^2\to\CI$ is
  continuous. So $y\mapsto RBR'(\cdot,y)$ is a smooth map into $\CI(X;E)$.
  Thus $B\in\Psido{-m}(X;F,E)$. Finally $\Tr(\pi)=\dim\Nul(A)$ and
  $\Tr(\pi')=\dim\Nul(A^*)=\dim\CI(X;F)/A\CI(X;E)$ by
  \cref{index:lem:proj-trace} and \eqref{index:eq:range}.
\end{proof}

Any two parametrices of $A$ differ by a smoothing operator. Indeed, if
$B_1A-\Id$ and $AB_2-\Id$ are smoothing then
$B_1-B_2=B_1(\Id-AB_2)-(\Id-B_1A)B_2\in\Psido{-\infty}$. In particular the
generalized inverse is a parametrix in the sense of \eqref{index:eq:parametrix}.

\begin{theorem}\label{index:thm:trace-formula}
  Let $A\in\Psido{m}(X;E,F)$ be elliptic and let $B\in\Psido{-m}(X;F,E)$ be
  any parametrix, so that $\Id_E-BA$ and $\Id_F-AB$ are smoothing. Then
  \begin{equation}\label{index:eq:trace-formula}
    \ind(A)=\Tr(\Id_E-BA)-\Tr(\Id_F-AB).
  \end{equation}
\end{theorem}
\begin{proof}
  For the generalized inverse $\tilde B$ of
  \cref{index:prop:generalized-inverse} this is the identity
  $\ind(A)=\Tr(\pi)-\Tr(\pi')$. For a general parametrix,
  $C=B-\tilde B\in\Psido{-\infty}(X;F,E)$ and
  \begin{equation*}
    \bigl(\Tr(\Id-BA)-\Tr(\Id-AB)\bigr)-
    \bigl(\Tr(\Id-\tilde BA)-\Tr(\Id-A\tilde B)\bigr)=\Tr(AC)-\Tr(CA)=0
  \end{equation*}
  by \cref{index:lem:trace-commute}.
\end{proof}

\begin{question}
  In finite dimensions, how does \eqref{index:eq:trace-formula} reproduce the
  familiar formula for the index?
\end{question}
\begin{answer}
  Let $A:\C^n\to\C^m$ and let $B:\C^m\to\C^n$ be any linear map. Then
  $\Tr(\Id_n-BA)-\Tr(\Id_m-AB)=n-m$, because $\Tr(AB)=\Tr(BA)$. By rank--nullity,
  $n-m=\dim\Nul A-\dim(\C^m/A\C^n)$. In infinite dimensions $n-m$ becomes
  ``$\infty-\infty$''. The parametrix makes $\Id-BA$ and $\Id-AB$ smoothing, hence
  trace class, so the two infinite traces of the identity cancel before any trace
  is taken.
\end{answer}

\begin{remark}\label{index:rem:sign}
  The right side of \eqref{index:eq:trace-formula} is $\Tr([A,B])$ in a formal
  sense: $\Id-BA$ and $\Id-AB$ are individually trace class, but $AB$ and $BA$
  are not. The formula shows that the index measures how far the trace
  identity $\Tr(AB)=\Tr(BA)$ fails when the operators are not smoothing.
\end{remark}

\section{Homotopy invariance of the index}\label{index:sec:homotopy}

\Cref{index:thm:trace-formula} writes the index as a continuous function of the
operator, as long as the parametrix can be chosen continuously. A continuous
integer-valued function on an interval is constant. So the only work is to
construct parametrices that depend smoothly on a parameter.

\subsection{Smooth families}
Let $Z$ be a manifold, possibly with boundary or corners (we mostly use
$Z=[0,1]$ and $Z=[0,1]\times X'$ with $X'$ compact). A \emph{smooth family}
$A\in\CI(Z;\Psido{m}(X;E,F))$ is a family $A_z\in\Psido{m}(X;E,F)$ whose
kernels have the following properties. Away from $Z\times\Diag$ they are
smooth on $Z\times(X^2\setminus\Diag)$. In local coordinates and
trivializations near the diagonal they are of the form
\begin{equation*}
  \int e^{i(x-y)\cdot\xi}a(z,x,\xi)\,\dbar\xi+K(z,x,y),
\end{equation*}
where $K$ is smooth in all variables and $a$ is a classical symbol of order
$m$ depending smoothly on $z$. That is, every $z$-derivative of $a$ satisfies
the symbol estimates locally uniformly in $z$, and the homogeneous terms of
the expansion of $a$ are smooth in $z$. For $m=-\infty$ this simply means
$A\in\CI(Z\times X^2;\Hom(E,F)\otimes\Omega_R)$.

All the constructions of \cref{ch:conormal,ch:manifolds} (composition,
adjoints, quantization of symbols, asymptotic summation, parametrices) are
given by formulas that are continuous in the symbol seminorms, so they apply
to smooth families with the result again smooth. The only point needing a
comment is asymptotic summation (Borel's lemma). The cut-offs
$\chi(\epsilon_j\xi)$ can be chosen with $\epsilon_j$ decreasing so fast that
the estimates hold uniformly on an exhaustion of $Z$ by compact sets. We
record the consequences we need.

\begin{lemma}\label{index:lem:family-parametrix}
  Let $A\in\CI(Z;\Psido{m}(X;E,F))$ be a smooth family with $A_z$ elliptic for
  every $z\in Z$.
  \begin{enumerate}[label=(\alph*)]
  \item There is a smooth family $B\in\CI(Z;\Psido{-m}(X;F,E))$ with
    $\Id-B_zA_z$ and $\Id-A_zB_z$ in $\CI(Z;\Psido{-\infty})$.
  \item If $m=0$, $E=F$ and every $A_z$ is invertible, then
    $z\mapsto A_z^{-1}$ is a smooth family in $\CI(Z;\Psido0(X;E))$.
  \end{enumerate}
  Moreover, given $a\in\CI(Z\times S^*X;\hom(\pi^*E,\pi^*F))$ there is a smooth
  family $Q(a)\in\CI(Z;\Psido{m}(X;E,F))$ with $\sigma_m(Q(a)_z)=a_z$,
  depending linearly on $a$.
\end{lemma}
\begin{proof}
  The quantization $Q$ is obtained by extending $a$ homogeneously of degree $m$,
  cutting off near the zero section, quantizing in coordinate charts and
  summing with a partition of unity. Each step preserves smooth dependence on
  $z$. For (a), let $B'=Q(\sigma_m(A)^{-1})$ (with $-m$ in place of $m$). Then
  $\Id-B'A\in\CI(Z;\Psido{-1})$, and the left parametrix
  $B_L\sim\sum_{j\ge0}(\Id-B'A)^jB'$ is an asymptotic sum of smooth families,
  so can be chosen smooth in $z$. Similarly for a right parametrix $B_R$, and
  $B_L-B_R=B_L(\Id-AB_R)-(\Id-B_LA)B_R$ is a smooth family of smoothing
  operators.

  For (b), let $G_z=A_z^{-1}$. It lies in $\Psido0$ by
  \cref{index:prop:generalized-inverse} with $\pi=\pi'=0$. A smooth family in
  $\Psido0$ is a smooth family of bounded operators on $L^2$, because the
  $L^2$ norm of an operator of order $0$ is bounded by a symbol seminorm. So
  $z\mapsto G_z$ is smooth into the bounded operators, being the inverse of a
  smooth family of invertible bounded operators. With $B$ from (a) and
  $R=\Id-BA$, $R'=\Id-AB$, the argument in the proof of
  \cref{index:prop:generalized-inverse} (with $\pi=\pi'=0$) gives
  \begin{equation*}
    G=G(AB+R')=B+GR'=B+(BA+R)GR'=B+BR'+RGR'.
  \end{equation*}
  Since $R_z$, $R'_z$ are smooth families of smoothing operators and $G_z$ is
  smooth in operator norm, the kernel of $RGR'$ is smooth in all variables.
  Hence $G$ is a smooth family.
\end{proof}

\subsection{Invariance}
\begin{theorem}\label{index:thm:homotopy}
  If $A\in\CI([0,1];\Psido{m}(X;E,F))$ is a smooth family with $A_t$ elliptic
  for every $t\in[0,1]$, then $\ind(A_t)$ is independent of $t$.
\end{theorem}
\begin{proof}
  Choose a smooth family of parametrices $B_t$ as in
  \cref{index:lem:family-parametrix}. By \cref{index:thm:trace-formula}
  \begin{equation*}
    \ind(A_t)=\Tr(\Id-B_tA_t)-\Tr(\Id-A_tB_t).
  \end{equation*}
  The arguments of the traces are continuous (indeed smooth) maps
  $[0,1]\to\Psido{-\infty}$ and the trace is continuous, so $\ind(A_t)$ is a
  continuous integer-valued function on $[0,1]$. Hence it is constant.
\end{proof}

\begin{discussion}
  Where was ellipticity used? Only in producing the smooth family of
  parametrices $B_t$. If ellipticity fails at one time $t_0$, then $\Id-B_tA_t$
  cannot be smoothing near $t_0$, the traces are undefined, and the index can
  jump. \Cref{index:fig:winding} shows the typical picture for the Toeplitz
  operators of \cref{index:sec:toeplitz}: along the family of symbols
  $a_t(\theta)=(1-t)+te^{i\theta}$ the index is $0$ for $t<\frac12$ and $-1$ for
  $t>\frac12$, and at $t=\frac12$ the symbol vanishes at one point.
\end{discussion}

\begin{figure}[ht]
\centering
\begin{tikzpicture}[scale=1.3,>=Stealth]
\draw[->,gray] (-1.3,0)--(1.5,0) node[right]{\small$\Re$};
\draw[->,gray] (0,-1.2)--(0,1.2) node[above]{\small$\Im$};
\fill (0,0) circle (1.2pt) node[below left]{\small$0$};
\foreach \t/\s in {0.2/densely dotted,0.35/densely dotted}{
  \draw[thick,\s] ({1-\t},0) circle (\t);}
\draw[very thick] (0.5,0) circle (0.5);
\fill[white] (0,0) circle (1.6pt); \draw (0,0) circle (1.6pt);
\foreach \t in {0.7,0.9}{\draw[thick,dashed] ({1-\t},0) circle (\t);}
\fill (1,0) circle (1.2pt) node[below right]{\small$1$};
\node[font=\footnotesize,align=left,anchor=west] at (2.2,0.9)
  {dotted: $t<\frac12$, $\wind=0$, $\ind_{\mathcal T}=0$};
\node[font=\footnotesize,align=left,anchor=west] at (2.2,0.45)
  {solid: $t=\frac12$, passes through $0$};
\node[font=\footnotesize,align=left,anchor=west] at (2.2,0)
  {dashed: $t>\frac12$, $\wind=1$, $\ind_{\mathcal T}=-1$};
\end{tikzpicture}
\caption{The images in $\C$ of the symbols $a_t(\theta)=(1-t)+te^{i\theta}$,
circles of centre $1-t$ and radius $t$. As long as the curve avoids $0$ (the
symbol is elliptic) its winding number about $0$, and with it the index of
$T_{a_t}$ (\cref{index:thm:toeplitz-index}), cannot change. The only way to
change the index is to pass through a non-elliptic symbol, here at
$t=\frac12$.}
\label{index:fig:winding}
\end{figure}

\begin{remark}\label{index:rem:continuous}
  The same conclusion holds for families that are only continuous in $t$ (in
  the topology of $\Psido m$). One way to see this is to note that such a
  family is norm-continuous as a family of Fredholm operators
  $\Hs m\to L^2$. Its index there equals the index on $\CI$, by the proof of
  \cref{index:prop:generalized-inverse}. The index of Fredholm operators is
  locally constant in the norm topology (\cref{ch:hilbert}).
\end{remark}

Two elliptic operators $A,A'\in\Psido m(X;E,F)$ with the same principal symbol
are joined by the smooth family $(1-t)A+tA'$, which has constant principal
symbol and so is elliptic. Hence $\ind(A)=\ind(A')$. More precisely, we have
the following.

\begin{corollary}\label{index:cor:symbol}
  Let $A_0,A_1\in\Psido m(X;E,F)$ be elliptic. Suppose that
  \begin{equation*}
    a\in\CI([0,1]\times S^*X;\hom(\pi^*E,\pi^*F))
  \end{equation*}
  is invertible at every point, with $a_0=\sigma_m(A_0)$ and $a_1=\sigma_m(A_1)$. Then
  $\ind(A_0)=\ind(A_1)$. Thus the index of an elliptic operator depends only
  on the homotopy class of its principal symbol in
  $\CI(S^*X;\Iso(\pi^*E,\pi^*F))$.
\end{corollary}
\begin{proof}
  With $Q$ as in \cref{index:lem:family-parametrix},
  \begin{equation*}
    A_t=Q(a)_t+(1-t)\bigl(A_0-Q(a)_0\bigr)+t\bigl(A_1-Q(a)_1\bigr)
  \end{equation*}
  is a smooth family with $\sigma_m(A_t)=a_t$ invertible, and $A_0$, $A_1$ are
  its end points. Apply \cref{index:thm:homotopy}.
\end{proof}

\begin{remark}\label{index:rem:order}
  (1) A continuous homotopy of invertible symbols between smooth end points can
  be approximated, with the same end points, by a smooth one, since the
  invertible elements form an open set. So continuous homotopies suffice in
  \cref{index:cor:symbol}.

  (2) The order is also irrelevant. For each $s\in\R$ there is an invertible
  elliptic $\Lambda_s\in\Psido{s}(X;E)$. Take any elliptic
  $\Lambda\in\Psido{s/2}(X;E)$ and set $\Lambda_s=\Lambda^*\Lambda+\pi$, with
  $\pi$ the orthogonal projection onto $\Nul(\Lambda)$. If
  $\Lambda_su=0$ then $\norm{\Lambda u}^2+\norm{\pi u}^2=0$, so $u=0$. Since
  $\Lambda_s$ is self-adjoint, \eqref{index:eq:range} shows it is also
  surjective. For $A\in\Psido m(X;E,F)$ elliptic,
  $A\Lambda_{-m}\in\Psido0(X;E,F)$ has the same index, since
  $\Nul(A\Lambda_{-m})=\Lambda_{-m}^{-1}\Nul(A)$ and the ranges coincide. So
  the index is a function on the homotopy classes of
  $\CI(S^*X;\Iso(\pi^*E,\pi^*F))$. Even this contains more information than
  the index sees. The index theorem eliminates the extra data by passing to
  K-theory, which we now introduce.
\end{remark}

\section{Smoothing groups and K-theory}\label{index:sec:K}

\Cref{index:cor:symbol} reduces the index to a function on homotopy classes of
invertible symbols. To compute it we need a good way of organizing such classes.
K-theory does this by ignoring the size of the matrices: a map $g$ into
$\GL(N,\C)$ is identified with its \emph{stabilization}
$g\oplus\Id\in\GL(N+1,\C)$. The group $\Ginf(Y;E)$ of invertible smoothing
perturbations of the identity contains all the groups $\GL(N,\C)$, nested inside
each other, so it performs this stabilization automatically.

\subsection{Odd K-theory}
Let $Y$ be a compact manifold of positive dimension and $E$ a vector bundle
over $Y$. Let $X$ be a compact manifold, possibly with boundary or corners. A
\emph{smooth map} $f:X\to\Ginf(Y;E)$ is a map such that
$f-\Id\in\CI(X;\Psido{-\infty}(Y;E))=\CI(X\times Y^2;\Hom(E)\otimes\Omega_R)$
and $f(x)$ is invertible for each $x$. Then $x\mapsto f(x)^{-1}-\Id$ is again
smooth. If $f=\Id+a$ and $f^{-1}=\Id+b$, the identity $b=-a+a^2+aba$ shows
that $b$ is smooth into $\Psido{-\infty}$ once it is continuous in $L^2$
norm, and the latter holds because inversion is continuous on bounded
operators. A smooth homotopy is a smooth map on $[0,1]\times X$. As in
\cref{ch:manifolds} we set
\begin{equation}\label{index:eq:K1}
  K^{-1}(X)=[X;\Ginf(Y;E)],
\end{equation}
the set of smooth homotopy classes of smooth maps. We recall the result of
\cref{ch:manifolds}: \eqref{index:eq:K1} is independent of $(Y,E)$. The
natural bijections are induced by finite rank approximation, which deforms
any smooth $f$ to a map of the form
\begin{equation}\label{index:eq:finite-rank}
  f(x)=\Id+\sum_{i,j=1}^N\bigl(g_{ij}(x)-\delta_{ij}\bigr)\phi_i\otimes\phi_j^*,
  \quad g\in\CI(X;\GL(N,\C)),
\end{equation}
for an orthonormal family $\phi_1,\dots,\phi_N\in\CI(Y;E)$, where
$\phi_i\otimes\phi_j^*$ denotes $u\mapsto\inner{u}{\phi_j}\phi_i$. Such an
$f$ acts as the matrix $g(x)$ on $W_N=\operatorname{span}\{\phi_i\}$ and as
the identity on $W_N^\perp$. It is invertible precisely when $g(x)$ is. To
see that a smooth $f=\Id+a$ is homotopic to a map of this form, approximate
$a$ uniformly on $X$ in $\CI(X;\Psido{-\infty})$ by
$a^{(N)}=\sum_{i,j\le N}c_{ij}(x)\phi_i\otimes\phi_j^*$, as in (iv). Since
$\Ginf$ is open and $X$ compact, $\Id+(1-s)a+sa^{(N)}$ is invertible for
$s\in[0,1]$ once $N$ is large.

\begin{proposition}\label{index:prop:K1-group}
  Pointwise multiplication induces a product on $K^{-1}(X)$ which makes it an
  abelian group, with identity the class of the constant map $\Id$ and
  inverse $[f]^{-1}=[f^{-1}]$. The product is also induced by the
  ``direct sum'': if $f$ and $f'$ are of the form \eqref{index:eq:finite-rank}
  with matrices $g$, $g'$ acting on orthogonal spaces $W$, $W'$, then
  $[f][f']$ is represented by the map acting as $g\oplus g'$ on $W\oplus W'$
  and as the identity on $(W\oplus W')^\perp$.
\end{proposition}
\begin{proof}
  If $f_s$, $f'_s$ are homotopies then so is $f_sf'_s$. Hence the product is
  well defined on homotopy classes. Associativity, the identity and the inverse
  are inherited from $\Ginf$. For commutativity, deform $f$ and $f'$ to the
  form \eqref{index:eq:finite-rank} with matrices $g$, $g'$ with respect to the
  same orthonormal family $\phi_1,\dots,\phi_N$. Enlarge the family to
  $\phi_1,\dots,\phi_{2N}$ (possible since $\dim Y>0$, so $\CI(Y;E)$ is
  infinite dimensional). Let $W_1$ and $W_2$ be spanned by the first and last
  $N$ elements. For $\theta\in[0,\pi/2]$ let $R_\theta$ be the unitary
  operator acting on $W_1\oplus W_2=\C^N\oplus\C^N$ by the block rotation
  \begin{equation*}
    R_\theta=\begin{pmatrix}\cos\theta&-\sin\theta\\
      \sin\theta&\cos\theta\end{pmatrix}
  \end{equation*}
  and as the identity on $(W_1\oplus W_2)^\perp$. Then $R_\theta-\Id$ is a
  finite rank smoothing operator and $R_\theta f'R_\theta^{-1}$ is a smooth
  homotopy in $\Ginf$. At $\theta=\pi/2$ it acts as $\Id_{W_1}\oplus g'$, with
  $g'$ now acting on $W_2$; this is a direct computation with
  $R_{\pi/2}=\left(\begin{smallmatrix}0&-1\\1&0\end{smallmatrix}\right)$.
  Thus $f'$ is homotopic to $f''=\Id_{W_1}\oplus g'$, which commutes with
  $f=g\oplus\Id_{W_2}$, and
  \begin{equation*}
    [f][f']=[ff'']=[g\oplus g']=[f''f]=[f'][f].
  \end{equation*}
  The middle term is the direct sum description.
\end{proof}

\begin{question}
  The groups $\GL(N,\C)$ and $\Ginf(Y;E)$ are not commutative. Why is
  $K^{-1}(X)$?
\end{question}
\begin{answer}
  Because commutativity is only needed up to homotopy, and there is room to move.
  The rotation $R_\theta$ slides $g'$ into a block $W_2$ orthogonal to the block
  $W_1$ on which $g$ acts, and maps acting on orthogonal blocks commute. The room
  comes from $\dim\CI(Y;E)=\infty$, which is why $\dim Y>0$ is assumed. Inside a
  fixed $\GL(N,\C)$, with $g$ and $g'$ using all $N$ dimensions, there is no
  such room and the argument fails.
\end{answer}

\begin{remark}\label{index:rem:pi-k}
  In terms of homotopy groups, the assertion that \eqref{index:eq:K1} is odd
  K-theory, i.e.\ that $\Ginf$ is a classifying space for $K^{-1}$, amounts
  to Bott periodicity in the form
  \begin{equation}\label{index:eq:homotopy-groups}
    \pi_k(\Ginf(Y;E))=\begin{cases}0&k\text{ even},\\ \Z&k\text{ odd}.
    \end{cases}
  \end{equation}
  This is proved in \cref{ch:bott}. In the standard topological treatment
  $K^{-1}(X)$ is defined by stabilizing maps into $\GL(N,\C)$ as
  $N\to\infty$ \cite{Atiyah-K}. The group $\Ginf$ makes the stabilization
  automatic: by \eqref{index:eq:finite-rank} every map is stably a map into
  some $\GL(N,\C)$, and the rotation argument above is the usual proof that
  the stabilized product is commutative.
\end{remark}

\subsection{Even K-theory}
Odd K-theory is built from invertible maps. Even K-theory, $K^0$, is usually
built from vector bundles (\cref{ch:bott}). The bridge is the \emph{clutching
construction}: a loop $g:\Sphere^1\to\GL(N,\C)$ glues the trivial bundles
$\C^N$ over the two hemispheres of $\Sphere^2$ along the equator, and every
bundle on $\Sphere^2$ arises in this way (\cref{bott:fig:clutching},
\cref{bott:lem:sphere-clutching}). So loops of invertible maps are transition
functions of bundles, and they give a model of even K-theory.

Let $\Sphere^1=\{e^{i\theta}\}$ with base point $1$.

\begin{definition}\label{index:def:K2}
  $K^{-2}(X)$ is the set of smooth homotopy classes of \emph{pointed} smooth
  maps $f:X\times\Sphere^1\to\Ginf(Y;E)$, meaning $f(x,1)=\Id$ for all $x\in X$,
  where homotopies are also required to be pointed at every time.
\end{definition}

Pointwise product makes $K^{-2}(X)$ a group. To see that it is abelian, and
how it is related to $K^{-1}$, we compare it with $K^{-1}(X\times\Sphere^1)$.

\begin{lemma}\label{index:lem:K1-split}
  The map $\Phi:K^{-1}(X)\times K^{-2}(X)\to K^{-1}(X\times\Sphere^1)$,
  $\Phi([g],[h])=[gh]$, where $g$ is regarded as a map on $X\times\Sphere^1$
  constant in $\theta$, is an isomorphism of groups. Its inverse is
  \begin{equation}\label{index:eq:split}
    [f]\mapsto\bigl([f_1],[f_2]\bigr),\quad f_1(x)=f(x,1),\quad
    f_2(x,\theta)=f(x,1)^{-1}f(x,\theta).
  \end{equation}
  In particular $K^{-2}(X)$ is abelian and
  $K^{-1}(X\times\Sphere^1)\cong K^{-1}(X)\oplus K^{-2}(X)$.
\end{lemma}
\begin{proof}
  $\Phi$ is well defined, since homotopies of $g$ and $h$ multiply to a
  homotopy of $gh$. The map \eqref{index:eq:split} is also well defined. If
  $F$ is a homotopy from $f$ to $f'$, then $F(\cdot,1)$ is a homotopy from
  $f_1$ to $f'_1$ and $F(\cdot,1)^{-1}F$ is a pointed homotopy from $f_2$ to
  $f'_2$. The two maps are mutually inverse at the level of maps, since
  $f=f_1f_2$ and $(gh)_1=g$, $(gh)_2=h$. Finally $\Phi$ is a homomorphism
  because $K^{-1}(X\times\Sphere^1)$ is abelian by
  \cref{index:prop:K1-group}:
  \begin{equation*}
    \Phi([g][g'],[h][h'])=[g][g'][h][h']=[g][h][g'][h']
    =\Phi([g],[h])\Phi([g'],[h']).
  \end{equation*}
  Since $\Phi$ restricted to $\{1\}\times K^{-2}(X)$ is an injective
  homomorphism into an abelian group, $K^{-2}(X)$ is abelian.
\end{proof}

\subsection{The loop group}
Pointed maps may be further normalized to be \emph{flat} at the base point,
i.e.\ to differ from the identity by a map vanishing to infinite order at
$X\times\{1\}$. Such maps are naturally Schwartz functions on $X\times\R$.

\begin{definition}\label{index:def:loop-group}
  The (smooth, flat, pointed) loop group of $\Ginf(Y;E)$ is
  \begin{equation*}
    \Gloop(Y;E)=\bigl\{a\in\cS(\R;\Psido{-\infty}(Y;E));\ \Id+a(t)\in
    \Ginf(Y;E)\ \forall\,t\in\R\bigr\},
  \end{equation*}
  where
  $\cS(\R;\Psido{-\infty}(Y;E))=\cS(\R\times Y^2;\Hom(E)\otimes\Omega_R)$ is
  the space of kernels that are Schwartz in $t$, uniformly with all
  derivatives in $Y^2$. The product is $(\Id+a)(\Id+b)=\Id+a+b+ab$, computed
  pointwise in $t$.
\end{definition}

\begin{lemma}\label{index:lem:loop-group}
  $\Gloop(Y;E)$ is an open subset of $\cS(\R;\Psido{-\infty}(Y;E))$ and a
  topological group.
\end{lemma}
\begin{proof}
  If $a\in\Gloop$ then $\norm{a(t)}_{L^2}\to0$ as $\abs t\to\infty$. So
  $\sup_t\norm{(\Id+a(t))^{-1}}<\infty$ by continuity of inversion, and every
  $a'$ with $\sup_t\norm{a'(t)-a(t)}_{L^2}$ small also has $\Id+a'(t)$
  invertible for all $t$. The supremum is bounded by a Schwartz seminorm, so
  $\Gloop$ is open. Products of Schwartz families are Schwartz. For the
  inverse $\Id+b(t)$ we have $b=-a+a^2+aba$ pointwise. Here $b(t)$ is bounded
  on $L^2$ uniformly in $t$, as are its $t$-derivatives, by repeated
  differentiation of $\partial_tb=-(\Id+b)(\partial_ta)(\Id+b)$. The kernel of
  $aba$ is then Schwartz in $t$, because each factor $a$ contributes Schwartz
  decay and $\CI$ regularity in the outer variable. So $b\in\Gloop$, and
  composition and inversion are continuous for the same reasons.
\end{proof}

\begin{proposition}\label{index:prop:K2-loop}
  Let $\gamma(t)=\exp\bigl(i\pi(1+t/\ang t)\bigr)$, a diffeomorphism of $\R$
  onto $\Sphere^1\setminus\{1\}$. Pulling back by $\gamma$ gives a
  bijection
  \begin{equation*}
    K^{-2}(X)\cong[X;\Gloop(Y;E)]
  \end{equation*}
  onto the smooth homotopy classes of smooth maps $X\to\Gloop(Y;E)$.
\end{proposition}
\begin{proof}
  First, pulling back by $\gamma$ identifies the functions on $\Sphere^1$
  vanishing to infinite order at $1$ with $\cS(\R)$, with parameters. Indeed,
  near $t=\pm\infty$ the angle $\theta=\pi(1+t/\ang t)$ satisfies
  $\theta\sim\pi/(2t^2)$ and $2\pi-\theta\sim\pi/(2t^2)$. Moreover
  $d\theta/dt=\pi\ang t^{-3}$, so $\partial_t$ and $\partial_\theta$ differ by
  factors that are polynomially bounded in $t$ or in $\theta^{-1}$. Infinite
  order vanishing in $\theta$ is therefore equivalent to rapid decay of all
  $t$-derivatives. Hence smooth maps $X\to\Gloop$ correspond exactly to
  smooth pointed maps $X\times\Sphere^1\to\Ginf$ that are flat at
  $X\times\{1\}$, and the same holds for homotopies.

  It remains to show that every pointed map is homotopic, through pointed
  maps, to a flat one, and that two flat maps which are pointed homotopic are
  homotopic through flat maps. Let $f=\Id+F$ be pointed. Then $F(x,1)=0$, so
  by compactness there is $\epsilon>0$ with $\norm{F(x,e^{i\theta})}_{L^2}<1$
  for $\abs\theta\le2\epsilon$. Choose $\rho\in\CI(\Sphere^1)$ with
  $0\le\rho\le1$, $\rho=1$ for $\abs\theta\le\epsilon$ and $\rho=0$ for
  $\abs\theta\ge2\epsilon$. Then
  \begin{equation*}
    f_s=\Id+(1-s\rho)F,\quad s\in[0,1],
  \end{equation*}
  is a pointed homotopy. Where $\rho\ne0$ we have
  $\norm{(1-s\rho)F}_{L^2}<1$, so $f_s$ is invertible. The end point $f_1$
  equals $\Id$ near $X\times\{1\}$, so it is flat. If $f$ is already flat,
  every $f_s$ is flat. Applying the same construction to a pointed homotopy
  between flat maps, with $X$ replaced by $[0,1]\times X$, gives a homotopy
  through flat maps from $f_1$ to $f'_1$. Combined with the flat homotopies
  $f_s$ and $f'_s$ at the ends, this connects $f$ to $f'$ through flat maps.
\end{proof}

\begin{remark}\label{index:rem:classifying}
  Thus, essentially by definition, $\Ginf(Y;E)$ and $\Gloop(Y;E)$ are
  classifying spaces for odd and even K-theory respectively. The second is the
  space of based loops in the first, so
  $\pi_k(\Gloop)=\pi_{k+1}(\Ginf)$, and \eqref{index:eq:homotopy-groups}
  becomes $\pi_0(\Gloop)=\Z$, $\pi_1(\Gloop)=0$ and so on. In \cref{ch:bott}
  $\Gloop$ is reinterpreted as the ``symbol group'' for elliptic Toeplitz
  operators on the circle with values in smoothing operators on $Y$. This
  leads to a \emph{classifying sequence}
  \begin{equation}\label{index:eq:classifying}
    \Ginf(\Sphere^1\times Y;E)\longrightarrow G\longrightarrow
    G^{-\infty}_{(1),0}(Y;E),
  \end{equation}
  in which $G$ is a weakly contractible group of invertible Toeplitz
  perturbations of the identity, and the subscript $0$ denotes the identity
  component of the loop group, on which the index vanishes. The last three
  sections of this chapter prepare for this.
\end{remark}

\section{Chern forms and the Fredholm determinant}\label{index:sec:chern}

For a loop $g$ in $\GL(N,\C)$ the basic invariant is the winding number of
$\det g$, which equals $\frac1{2\pi i}\int\tr(g^{-1}dg)$. This section builds the
analogue on $\Ginf$, with the trace of smoothing operators in place of the matrix
trace. The forms $\Tr((A^{-1}dA)^{2k-1})$ give de~Rham classes on the K-groups,
and the first of them, integrated along paths, gives the Fredholm determinant.

\subsection{Forms with values in smoothing operators}
Since $\Ginf=\Ginf(Y;E)$ is an open subset of $\Id+\Psido{-\infty}(Y;E)$, we
may treat it as a (Fr\'echet) manifold whose tangent space at every point is
$\Psido{-\infty}(Y;E)$. The tangent vector to a smooth curve $\Id+a_t$ at
$A=\Id+a_0$ is $dA=\frac{d}{dt}a_t|_{t=0}$. The \emph{Maurer--Cartan form}
$A^{-1}dA$, the map $T_A\Ginf=\Psido{-\infty}\ni b\mapsto A^{-1}b$, is a
$\Psido{-\infty}$-valued $1$-form. It is invariant under left
multiplication: if $L_B(A)=BA$ then $L_B^*(A^{-1}dA)=A^{-1}B^{-1}B\,dA=A^{-1}dA$.
Products of such forms are defined by composition of the values combined with
antisymmetrization, and $\Tr$ is applied to the values. Explicitly, the
\emph{(unnormalized) odd Chern forms} are
\begin{equation}\label{index:eq:chern-forms}
  \begin{gathered}
  u_{2k-1}=\Tr\bigl((A^{-1}dA)^{2k-1}\bigr),\\
  u_{2k-1}(b_1,\dots,b_{2k-1})=\sum_{\sigma\in\Sigma_{2k-1}}\sgn(\sigma)
  \Tr\bigl(A^{-1}b_{\sigma(1)}A^{-1}b_{\sigma(2)}\cdots A^{-1}b_{\sigma(2k-1)}
  \bigr),
  \end{gathered}
\end{equation}
for $k=1,2,\dots$. They are smooth by the smoothness of composition, inversion
and the trace.

To avoid differential calculus on Fr\'echet manifolds we work with the
pull-backs of these forms under smooth maps $f:X\to\Ginf(Y;E)$, where $X$ is a
finite dimensional manifold. Let
\begin{equation*}
  \Lambda^p(X;\Psido{-\infty}(Y;E))=\CI\bigl(X\times Y^2;\Lambda^pT^*X\boxtimes
  (\Hom(E)\otimes\Omega_R)\bigr)
\end{equation*}
be the $p$-forms on $X$ with values in smoothing operators on $Y$. We also
allow $0$-forms of the form $c\Id+a$. The wedge product of
$\alpha\otimes a$ and $\beta\otimes b$ is $\alpha\wedge\beta\otimes ab$,
extended by linearity and locally finite sums. The differential $d$ acts in
the $X$ variables, and $\Tr$ acts on the values, producing ordinary forms. The
following are immediate from the corresponding identities for scalar forms
and from $\Tr(ab)=\Tr(ba)$:
\begin{equation}\label{index:eq:graded}
  \begin{aligned}
  d(\omega\wedge\eta)&=d\omega\wedge\eta+(-1)^p\omega\wedge d\eta,\qquad
  d\Tr(\omega)=\Tr(d\omega),\\
  \Tr(\omega\wedge\eta)&=(-1)^{pq}\Tr(\eta\wedge\omega),
  \end{aligned}
\end{equation}
for $\omega$ of degree $p$, $\eta$ of degree $q$, and at least one of them
smoothing-valued. For the last identity, write locally
$\omega=\sum\alpha_I\otimes a_I$ and $\eta=\sum\beta_J\otimes b_J$. Then
$\Tr(\omega\wedge\eta)=\sum\alpha_I\wedge\beta_J\Tr(a_Ib_J)
=(-1)^{pq}\sum\beta_J\wedge\alpha_I\Tr(b_Ja_I)$. With the antisymmetrization
convention of \eqref{index:eq:chern-forms}, the pull-back of $u_{2k-1}$ by $f$
is
\begin{equation*}
  f^*u_{2k-1}=\Tr\bigl(\theta^{2k-1}\bigr),\quad \theta=f^{-1}df\in
  \Lambda^1(X;\Psido{-\infty}(Y;E)).
\end{equation*}

\begin{lemma}\label{index:lem:closed}
  For every smooth $f:X\to\Ginf(Y;E)$, with $\theta=f^{-1}df$:
  \begin{enumerate}[label=(\alph*)]
  \item $d\theta=-\theta\wedge\theta$;
  \item $\Tr(\theta^{2k})=0$ for $k\ge1$;
  \item $f^*u_{2k-1}=\Tr(\theta^{2k-1})$ is closed;
  \item $f^*u_{2k-1}$ is unchanged if $f$ is replaced by $Bf$ or $fB$ with
    $B\in\Ginf(Y;E)$ fixed. That is, $u_{2k-1}$ is bi-invariant.
  \end{enumerate}
\end{lemma}
\begin{proof}
  Differentiating $f^{-1}f=\Id$ gives $d(f^{-1})=-f^{-1}(df)f^{-1}$, so
  $d\theta=d(f^{-1})\wedge df=-\theta\wedge\theta$. By \eqref{index:eq:graded}
  with $\omega=\theta$ and $\eta=\theta^{2k-1}$,
  $\Tr(\theta^{2k})=(-1)^{2k-1}\Tr(\theta^{2k})$, which is (b). For (c), the
  Leibniz rule and (a) give
  \begin{equation*}
    d(\theta^{n})=\sum_{j=0}^{n-1}(-1)^j\theta^j\wedge d\theta\wedge
    \theta^{n-1-j}=-\Bigl(\sum_{j=0}^{n-1}(-1)^j\Bigr)\theta^{n+1}.
  \end{equation*}
  For $n=2k-1$ odd the sum is $1$, so $d\Tr(\theta^{2k-1})=-\Tr(\theta^{2k})=0$
  by (b). For (d), $(Bf)^{-1}d(Bf)=\theta$, while
  $(fB)^{-1}d(fB)=B^{-1}\theta B$ and
  $\Tr((B^{-1}\theta B)^{2k-1})=\Tr(B^{-1}\theta^{2k-1}B)=\Tr(\theta^{2k-1})$.
\end{proof}

\begin{lemma}\label{index:lem:transgression}
  Let $F:[0,1]\times X\to\Ginf(Y;E)$ be smooth and $f_t=F(t,\cdot)$. Then
  \begin{equation*}
    f_1^*u_{2k-1}-f_0^*u_{2k-1}=d\int_0^1v_t\,dt,\quad
    v_t=(2k-1)\Tr\bigl(f_t^{-1}\partial_tf_t\,(f_t^{-1}d_Xf_t)^{2k-2}\bigr).
  \end{equation*}
\end{lemma}
\begin{proof}
  On $[0,1]\times X$ write $\Theta=F^{-1}dF=dt\,\alpha+\theta_t$ with
  $\alpha=f_t^{-1}\partial_tf_t$ and $\theta_t=f_t^{-1}d_Xf_t$. Expanding
  $\Tr(\Theta^{2k-1})$, the terms containing $dt$ are
  \begin{equation*}
    \sum_{j=0}^{2k-2}\Tr(\theta_t^j\wedge dt\,\alpha\,\theta_t^{2k-2-j})
    =dt\wedge\sum_j(-1)^j\Tr(\theta_t^j\alpha\theta_t^{2k-2-j})
    =dt\wedge v_t.
  \end{equation*}
  The last step uses \eqref{index:eq:graded} to move $\theta_t^j$ to the end
  at the cost of $(-1)^{j(2k-2-j)}=(-1)^j$. So
  $\Tr(\Theta^{2k-1})=dt\wedge v_t+f_t^*u_{2k-1}$, and this is closed by
  \cref{index:lem:closed}. Its $dt$-component gives
  $\partial_t(f_t^*u_{2k-1})-d_Xv_t=0$. Integrating in $t$ proves the
  lemma.
\end{proof}

\begin{proposition}\label{index:prop:odd-chern}
  For each $k\in\N$ the de~Rham class of $f^*u_{2k-1}$ defines a group
  homomorphism
  \begin{equation*}
    U_{2k-1}:K^{-1}(X)\longrightarrow H^{2k-1}_{\mathrm{dR}}(X;\C),\quad
    [f]\mapsto[f^*u_{2k-1}].
  \end{equation*}
\end{proposition}
\begin{proof}
  The class is well defined on homotopy classes by
  \cref{index:lem:transgression}. For additivity, deform $f$ and $f'$ as in the
  proof of \cref{index:prop:K1-group} so that $f=g\oplus\Id$ and
  $f'=\Id\oplus g'$ act on orthogonal finite dimensional spaces $W_1$, $W_2$.
  Then $\theta=f^{-1}df$ and $\theta'=(f')^{-1}df'$ satisfy
  $\theta\wedge\theta'=0=\theta'\wedge\theta$, and $(ff')^{-1}d(ff')=\theta+\theta'$.
  Hence
  $(\theta+\theta')^{2k-1}=\theta^{2k-1}+(\theta')^{2k-1}$ and
  $(ff')^*u_{2k-1}=f^*u_{2k-1}+(f')^*u_{2k-1}$.
\end{proof}

\begin{remark}\label{index:rem:odd-chern-character}
  For a map of the form \eqref{index:eq:finite-rank},
  $f^*u_{2k-1}=\tr((g^{-1}dg)^{2k-1})$ is the familiar form built from the
  matrix-valued map $g$, since the trace of a finite rank operator is the
  matrix trace. With suitable normalizing constants $c_k$ the combination
  $\sum_kc_kU_{2k-1}:K^{-1}(X)\to H^{\mathrm{odd}}_{\mathrm{dR}}(X;\C)$ is the
  \emph{odd Chern character}; the constants are fixed in \cref{ch:families}.
  It is a standard result of topology that its range spans
  $H^{\mathrm{odd}}(X;\C)$ and that its null space is the torsion subgroup of
  $K^{-1}(X)$ (see \cite{Atiyah-K}). We do not use this.
\end{remark}

\subsection{Even Chern forms}
Even forms are defined in the same way on the loop group. A smooth map
$f:X\to\Gloop(Y;E)$ is a map $F:X\times\R\to\Ginf(Y;E)$ with $F-\Id$
Schwartz in $t$, smooth in $x$. Set
\begin{equation}\label{index:eq:even-forms}
  f^*u_{2k}=\int_\R\Tr\bigl((F^{-1}d_XF)^{2k}\,F^{-1}\partial_tF\bigr)\,dt,
  \quad k=0,1,2,\dots.
\end{equation}
The integrand is Schwartz in $t$. By the computation in the proof of
\cref{index:lem:transgression}, with $(2k+1)$ in place of $(2k-1)$, we have
$\Tr((F^{-1}dF)^{2k+1})=dt\wedge v+w$ on $X\times\R$, where
$v=(2k+1)\Tr(F^{-1}\partial_tF(F^{-1}d_XF)^{2k})$. Since
$\Tr(\alpha\theta^{2k})=\Tr(\theta^{2k}\alpha)$, this gives
\begin{equation*}
  (2k+1)f^*u_{2k}=\int_\R v\,dt.
\end{equation*}
This is the integral over the fibres of $X\times\R\to X$ of the closed form
$F^*u_{2k+1}$. Because $\partial_tw=d_Xv$ and $w\to0$ rapidly as
$\abs t\to\infty$, we get $d_X\int_\R v\,dt=\int_\R\partial_tw\,dt=0$. So
$f^*u_{2k}$ is closed. The argument of \cref{index:lem:transgression},
applied on $[0,1]\times X\times\R$, shows that its class is a homotopy
invariant. Exactly as in \cref{index:prop:odd-chern}, using
\cref{index:prop:K2-loop}, we obtain homomorphisms
\begin{equation*}
  U_{2k}:K^{-2}(X)\longrightarrow H^{2k}_{\mathrm{dR}}(X;\C),\quad k\in\N_0.
\end{equation*}
A suitable combination of them is the (ordinary) Chern character.

The simplest and most fundamental cases are the first odd form and its even
counterpart:
\begin{equation}\label{index:eq:u1-u0}
  u_1=\Tr(A^{-1}dA)\ \text{on }\Ginf,\qquad
  u_0(A)=\int_\R\Tr\bigl(A(t)^{-1}\partial_tA(t)\bigr)\,dt,\ A\in\Gloop.
\end{equation}

\begin{proposition}\label{index:prop:integrality}
  For every piecewise smooth closed curve $\gamma:\Sphere^1\to\Ginf(Y;E)$,
  \begin{equation*}
    \int_\gamma u_1\in2\pi i\Z.
  \end{equation*}
  Consequently $u_0/2\pi i$ is an integer-valued, locally constant function
  on $\Gloop(Y;E)$.
\end{proposition}
\begin{proof}
  The integral is unchanged by reparametrization, so we may assume $\gamma$
  smooth (round off corners by a reparametrization that is flat at them). By
  \cref{index:lem:transgression} with $X=\Sphere^1$ and $k=1$, homotopic loops
  have equal integrals. By finite rank approximation we may therefore replace
  $\gamma$ by a loop of the form \eqref{index:eq:finite-rank} with
  $g:\Sphere^1\to\GL(N,\C)$. Then $\gamma^*u_1=\tr(g^{-1}dg)$. By the matrix
  identity $d\det g=\det g\,\tr(g^{-1}dg)$,
  \begin{equation*}
    \int_\gamma u_1=\int_{\Sphere^1}\frac{d\det g}{\det g}=2\pi i\,
    \wind(\det g),
  \end{equation*}
  where $\wind$ is the winding number about $0$. If $A\in\Gloop$, then
  $u_0(A)$ is the integral of $u_1$ over the closed loop $A\circ\gamma^{-1}$ in
  the notation of \cref{index:prop:K2-loop}, extended by $\Id$ at $1$, so it
  lies in $2\pi i\Z$. It is continuous on $\Gloop$, hence locally
  constant.
\end{proof}

\begin{remark}\label{index:rem:u0-index}
  We shall see in \cref{ch:bott} that $u_0/2\pi i:\pi_0(\Gloop)\to\Z$ is a
  bijection, so it labels the components of the loop group. It is the
  simplest case of the index formula: under the identification of
  \cref{index:rem:classifying} it computes the index of the corresponding
  Toeplitz operator, up to sign.
\end{remark}

\subsection{The Fredholm determinant}
\Cref{index:prop:integrality} yields a direct construction of the Fredholm
determinant on $\Ginf(Y;E)$. Any two points of $\Ginf$ are joined by a
piecewise linear path. Indeed $\Ginf$ is connected, and each point
$A\in\Ginf$ has a convex neighbourhood
$\{A+b;\ \norm b_{L^2}<\norm{A^{-1}}^{-1}\}$ inside $\Ginf$. After
reparametrization such a path is smooth.

\begin{discussion}
  In finite dimensions $d\log\det A=\tr(A^{-1}dA)$, so $\log\det A$ is the
  integral of $\tr(A^{-1}dA)$ along a path from $\Id$ to $A$. The logarithm is
  multivalued: two paths differ by a loop, and the integral over a loop is
  $2\pi i$ times a winding number. The exponential does not see this ambiguity.
  The same argument works on $\Ginf$ because \cref{index:prop:integrality} says
  that the periods of $u_1$ lie in $2\pi i\Z$. We never need a trace of
  $\log A$ or a product of eigenvalues.
\end{discussion}

\begin{theorem}\label{index:thm:det}
  For $A\in\Ginf(Y;E)$ let $\gamma:[0,1]\to\Ginf(Y;E)$ be any piecewise smooth
  path with $\gamma(0)=\Id$ and $\gamma(1)=A$, and set
  \begin{equation}\label{index:eq:det}
    \det(A)=\exp\Bigl(\int_\gamma u_1\Bigr)=\exp\Bigl(\int_0^1
    \Tr\bigl(\gamma(t)^{-1}\gamma'(t)\bigr)\,dt\Bigr).
  \end{equation}
  This is independent of $\gamma$. It defines a smooth homomorphism
  $\det:\Ginf(Y;E)\to\C^*$ with $d\log\det=u_1$, i.e.\
  $\frac{d}{ds}\det(A_s)=\det(A_s)\Tr(A_s^{-1}\frac{dA_s}{ds})$ along smooth
  curves. If $A=\Id+a$ with $a$ of finite rank and range in a finite
  dimensional space $W\subset\CI(Y;E)$, then $\det(A)$ is the determinant of
  the restriction $A|_W:W\to W$.
\end{theorem}
\begin{proof}
  If $\gamma_1$, $\gamma_2$ are two such paths, then $\gamma_1$ followed by
  $\gamma_2$ reversed is a closed piecewise smooth loop. By
  \cref{index:prop:integrality} the two integrals differ by an element of
  $2\pi i\Z$. For multiplicativity, if $\gamma_A$ and $\gamma_B$ join $\Id$ to
  $A$ and $B$, then $\gamma_A\gamma_B$ joins $\Id$ to $AB$ and
  \begin{equation*}
    (\gamma_A\gamma_B)^{-1}(\gamma_A\gamma_B)'=\gamma_B^{-1}
    (\gamma_A^{-1}\gamma_A')\gamma_B+\gamma_B^{-1}\gamma_B'.
  \end{equation*}
  Taking traces gives $\det(AB)=\det(A)\det(B)$. For smoothness near
  $A_0\in\Ginf$, choose the path to pass through $A_0$ and then follow the
  segment to $A$. This gives
  \begin{equation*}
    \det(A)=\det(A_0)\exp\Bigl(\int_0^1\Tr\bigl((A_0+s(A-A_0))^{-1}(A-A_0)
    \bigr)\,ds\Bigr),
  \end{equation*}
  which is smooth in $A$ near $A_0$. Differentiating along a curve gives
  $d\log\det=u_1$. Finally, suppose $a$ has range in $W$. Then $A$ preserves
  $W$ and induces the identity on $\CI(Y;E)/W$. In terms of a decomposition
  $\CI=W\oplus W'$ it is upper triangular,
  $\left(\begin{smallmatrix}g&c\\0&\Id\end{smallmatrix}\right)$ with
  $g=A|_W\in\GL(W)$. Since $\GL(W)$ is connected we can join $A$ to $\Id$ by
  $\gamma(t)=\left(\begin{smallmatrix}g_t&tc\\0&\Id\end{smallmatrix}\right)$
  with $g_t$ a path in $\GL(W)$ from $\Id$ to $g$. Then
  $\gamma^{-1}\gamma'$ has range in $W$, and its trace is
  $\tr(g_t^{-1}g_t')$. Integrating gives $\det(A)=\det(g)$.
\end{proof}

By \cref{index:thm:det} the determinant agrees with the matrix determinant
under any finite rank embedding $\GL(N,\C)\hookrightarrow\Ginf$. It is the
same function as the Fredholm determinant constructed for the residual
isotropic algebra in \cref{ch:determinants}; the characterizing properties
are identical. It extends to all smoothing perturbations of the identity,
vanishing exactly on the non-invertible ones.

\begin{theorem}\label{index:thm:det-extension}
  The function $a\mapsto\det(\Id+a)$, defined on
  $\{a\in\Psido{-\infty}(Y;E);\ \Id+a\in\Ginf\}$, extends uniquely to a
  continuous function on $\Psido{-\infty}(Y;E)$. The extension is smooth,
  entire along complex lines (i.e.\ $z\mapsto\det(\Id+a+zb)$ is entire), and
  vanishes exactly where $\Id+a$ is not invertible. Moreover $\Ginf$ is dense
  in $\Id+\Psido{-\infty}(Y;E)$, and
  $\frac{d}{ds}\det(\Id+sa)|_{s=0}=\Tr(a)$.
\end{theorem}
\begin{proof}
  Fix $a_0\in\Psido{-\infty}(Y;E)$. Using (iv) write $a_0=p_0+r_0$ with $p_0$
  of finite rank, with range in a finite dimensional $W\subset\CI(Y;E)$, and
  $\norm{r_0}_{L^2}<1/2$. For $a$ in the neighbourhood
  $\mathcal U=\{\norm{a-a_0}_{L^2}<1/2\}$ put $r=a-p_0$. Then $\norm r<1$, so
  $\Id+r\in\Ginf$, and
  \begin{equation*}
    \Id+a=\bigl(\Id+p_0(\Id+r)^{-1}\bigr)(\Id+r).
  \end{equation*}
  The first factor is $\Id$ plus an operator with range in $W$. As in the
  proof of \cref{index:thm:det} it is invertible exactly when its restriction
  $g(a)$ to $W$ is. Define on $\mathcal U$
  \begin{equation*}
    D(a)=\det\bigl(g(a)\bigr)\det(\Id+r),\quad
    g(a)=\bigl(\Id+p_0(\Id+a-p_0)^{-1}\bigr)\big|_W.
  \end{equation*}
  This is smooth and holomorphic along complex lines, because
  $a\mapsto(\Id+a-p_0)^{-1}$ is. It vanishes exactly where $\Id+a$ is not
  invertible. Where $\Id+a\in\Ginf$ it equals $\det(\Id+a)$ by
  \cref{index:thm:det}, which gives both multiplicativity and the value on
  finite rank perturbations.

  Next we show that $\{a\in\mathcal U;\ \Id+a\in\Ginf\}$ is dense in $\mathcal
  U$. Let $\pi_W$ be the orthogonal projection onto $W$, a finite rank
  smoothing operator, and for $a\in\mathcal U$ set
  $a_\epsilon=a+\epsilon\pi_W(\Id+r)$, $r=a-p_0$. Then
  \begin{equation*}
    \Id+a_\epsilon=\bigl(\Id+p_0(\Id+r)^{-1}+\epsilon\pi_W\bigr)(\Id+r),
  \end{equation*}
  and the first factor is $\Id$ plus an operator with range in $W$, whose
  restriction to $W$ is $g(a)+\epsilon\Id_W$. This is invertible for all
  small $\epsilon\ne0$, and $a_\epsilon\to a$.
  Two continuous functions agreeing on a dense set agree. So the local
  definitions agree on overlaps and define a global extension, which is
  unique for the same reason. Finally
  $\det(\Id+sa)=\exp\int_0^s\Tr((\Id+\sigma a)^{-1}a)\,d\sigma$ for small
  $s$, whose derivative at $0$ is $\Tr(a)$.
\end{proof}

\section{The Szeg\H{o} projector and Toeplitz operators}\label{index:sec:toeplitz}

\subsection{Motivation}
The next step towards the index theorem is a group-level version of the symbol
sequence
\begin{equation*}
  \Psido{-1}(Z;\C^N)\longrightarrow\Psido0(Z;\C^N)\xrightarrow{\ \sigma_0\ }
  \CI(S^*Z;M(N,\C)).
\end{equation*}
Passing to invertible elements gives
\begin{equation}\label{index:eq:group-sequence}
  G^{-1}(Z;\C^N)\longrightarrow G^0(Z;\C^N)\longrightarrow
  \CI(S^*Z;\GL(N,\C)),
\end{equation}
where $G^0(Z;\C^N)$ is the group of invertible elliptic elements of
$\Psido0(Z;\C^N)$ (their inverses lie in $\Psido0$ by
\cref{index:prop:generalized-inverse}) and
$G^{-1}(Z;\C^N)=G^0(Z;\C^N)\cap(\Id+\Psido{-1}(Z;\C^N))$. The second map is
in general not surjective. An elliptic symbol is the symbol of an
invertible operator if and only if the index of some (hence every) operator
with that symbol vanishes. So the index is the obstruction to exactness of
\eqref{index:eq:group-sequence}.

We would like to arrange that the middle group is contractible and the group on
the right is essentially a loop group $\Gloop$. This requires two things.
First, $S^*Z$ should be a circle. Second, $\C^N$ should be replaced by an
infinite dimensional space, so that $\GL(N,\C)$ becomes one of the groups
$\Ginf$ (this is done in \cref{ch:bott}). A compact manifold with
$S^*Z=\Sphere^1$ does not exist, since $Z$ would be a compact
$1$-dimensional manifold and $S^*\Sphere^1=\Sphere^1\sqcup\Sphere^1$. There
are two ways round this:
\begin{enumerate}[label=(\Alph*)]
\item replace the circle by the line $\R$ and use the isotropic calculus of
  \cref{ch:isotropic}, for which the ``cosphere'' is the boundary circle of
  the radial compactification of $T^*\R=\R^2$;
\item keep the circle but discard half of $S^*\Sphere^1$, which leads to the
  Toeplitz algebra.
\end{enumerate}
We follow (B). It is less transparent geometrically but has a long history. A
third route uses the Heisenberg calculus \cite{Epstein-Melrose}.

\subsection{The Hardy space and the Szeg\H{o} projector}
Expanding in Fourier series, $\CI(\Sphere^1)$ is the direct sum of
\begin{equation*}
  \Hardy(\Sphere^1)=\Bigl\{\sum_{k\ge0}c_ke^{ik\theta}\Bigr\},\qquad
  \Hardym(\Sphere^1)=\Bigl\{\sum_{k<0}c_ke^{ik\theta}\Bigr\},
\end{equation*}
the coefficients being rapidly decreasing. The \emph{Szeg\H{o} projector} is
\begin{equation}\label{index:eq:szego}
  S:\CI(\Sphere^1)\longrightarrow\Hardy(\Sphere^1),\quad
  S\Bigl(\sum_{k\in\Z}c_ke^{ik\theta}\Bigr)=\sum_{k\ge0}c_ke^{ik\theta}.
\end{equation}
It satisfies $S^2=S=S^*$ (with respect to $d\theta$), its null space is
$\Hardym$ and its range is $\Hardy$. The product of two elements of $\Hardy$ is
again in $\Hardy$, so the Hardy space $\Hardy(\Sphere^1)$ is a ring. It
consists of the boundary values of holomorphic functions on the unit disc
$\mathbb D\subset\C$ (here $\mathbb{D}$ is the open disc and
$\overline{\mathbb D}$ its closure), so it should be thought of as the
functions on ``half'' of $\Sphere^1$. The cotangent bundle is
$T^*\Sphere^1=\Sphere^1\times\R_\xi$, and the Fourier variable $k$
corresponds to $\xi$. This suggests that $S$ is microlocally the identity on
the positive half $S^*_+\Sphere^1=\{\xi>0\}$ and zero on the negative half
$S^*_-\Sphere^1$ (\cref{index:fig:hardy}).

\begin{figure}[ht]
\centering
\begin{tikzpicture}[>=Stealth]
\begin{scope}
\draw[->,gray] (-3.3,0)--(3.4,0) node[right]{\small$k$};
\foreach \k in {-3,...,-1}{\draw (\k*0.9,0) circle (2.2pt);}
\foreach \k in {0,...,3}{\fill (\k*0.9,0) circle (2.2pt);}
\foreach \k in {-3,...,3}{\node[below,font=\footnotesize] at (\k*0.9,-0.12){$\k$};}
\draw[decorate,decoration={brace,amplitude=4pt}] (-0.1,0.25)--(2.8,0.25)
  node[midway,above=4pt,font=\footnotesize]{$\Hardy$: kept by $S$};
\draw[decorate,decoration={brace,amplitude=4pt}] (-2.8,0.25)--(-0.8,0.25)
  node[midway,above=4pt,font=\footnotesize]{$\Hardym$: killed};
\node[font=\small] at (0,-1.1) {(a) Fourier modes $e^{ik\theta}$};
\end{scope}
\begin{scope}[shift={(5.6,0.1)}]
\draw[gray] (0,0) circle (0.9);
\draw[very thick] (0,0) circle (1.45);
\draw[thick,densely dashed] (0,0) circle (0.4);
\foreach \a in {30,90,...,330}{
  \draw[->] (\a:0.9)--(\a:1.4);
  \draw[->,densely dashed] (\a:0.9)--(\a:0.45);}
\draw (40:1.45)--(1.9,1.2) node[right,font=\footnotesize]{$S^*_+$: $\sigma_0(S)=1$};
\draw (0:0.9)--(1.9,0) node[right,font=\footnotesize]{$\Sphere^1$ (zero section)};
\draw (-60:0.4)--(1.9,-1.2) node[right,font=\footnotesize]{$S^*_-$: $\sigma_0(S)=0$};
\node[font=\small] at (0.9,-1.95) {(b) $T^*\Sphere^1$, $\xi>0$ outward};
\end{scope}
\end{tikzpicture}
\caption{The Szeg\H{o} projector. (a) On Fourier series, $S$ keeps the modes
$k\ge0$ and kills the modes $k<0$. (b) The cotangent bundle
$T^*\Sphere^1=\Sphere^1\times\R_\xi$, drawn with $\xi>0$ pointing outwards. The
cosphere bundle has two components, $S^*_+$ (outer circle) and $S^*_-$ (inner
circle); the principal symbol of $S$ is $1$ on the first and $0$ on the second
(\cref{index:thm:szego}).}
\label{index:fig:hardy}
\end{figure}

\begin{proposition}\label{index:prop:hardy}
  Restriction to the boundary is an isomorphism
  \begin{equation*}
    \CI_{\mathrm{hol}}(\overline{\mathbb D})=\{u\in\CI(\overline{\mathbb D});\
    \bar\partial u=0\text{ in }\mathbb D\}\longrightarrow\Hardy(\Sphere^1).
  \end{equation*}
  For $f\in\CI(\Sphere^1)$ the Cauchy integral
  \begin{equation*}
    \tilde S(f)(z)=\frac1{2\pi i}\int_{\abs\tau=1}\frac{f(\tau)}{\tau-z}\,d\tau,
    \quad\abs z<1,
  \end{equation*}
  is the holomorphic extension of $Sf$, and
  $Sf(e^{i\theta})=\lim_{r\uparrow1}\tilde Sf(re^{i\theta})$.
\end{proposition}
\begin{proof}
  If $a=\sum_{k\ge0}a_ke^{ik\theta}\in\Hardy$ then the coefficients are
  rapidly decreasing. Hence $u_a(z)=\sum_{k\ge0}a_kz^k$ converges on
  $\overline{\mathbb D}$ with all derivatives, is holomorphic in $\mathbb D$,
  and restricts to $a$. Conversely, if $u\in\CI_{\mathrm{hol}}(\overline{\mathbb
  D})$ has the Taylor series $\sum_{k\ge0}b_kz^k$ in $\mathbb D$, then for
  $r<1$ and $k\in\Z$
  \begin{equation*}
    \frac1{2\pi}\int_0^{2\pi}e^{-ik\theta}u(re^{i\theta})\,d\theta=
    \begin{cases}b_kr^k&k\ge0,\\0&k<0.\end{cases}
  \end{equation*}
  Letting $r\uparrow1$, the Fourier coefficients of $u|_{\Sphere^1}$ are $b_k$
  for $k\ge0$ and $0$ for $k<0$. So the boundary value lies in $\Hardy$, and
  it vanishes only if $u=0$. Thus restriction is an isomorphism. For the
  Cauchy integral, by rapid convergence of the Fourier series we can integrate
  term by term. For $\abs z<1$,
  \begin{equation*}
    \frac1{2\pi i}\int_{\abs\tau=1}\frac{\tau^k}{\tau-z}\,d\tau=
    \begin{cases}z^k&k\ge0,\\0&k<0,\end{cases}
  \end{equation*}
  by the residue theorem. For $k<0$ the integrand $\tau^k/(\tau-z)$ is
  holomorphic outside the disc and $O(\abs\tau^{-2})$ at infinity, so we
  deform the contour to infinity. So $\tilde S(f)=u_{Sf}$, and the boundary
  values follow from the first part.
\end{proof}

\begin{theorem}\label{index:thm:szego}
  $S\in\Psido0(\Sphere^1)$, and its principal symbol is
  \begin{equation*}
    \sigma_0(S)=\begin{cases}1&\text{on }S^*_+\Sphere^1=\{\xi>0\},\\
    0&\text{on }S^*_-\Sphere^1=\{\xi<0\}.\end{cases}
  \end{equation*}
  More precisely, near the diagonal the kernel of $S$ is
  $\int e^{i(\theta-\theta')\xi}\chi(\xi)\,\dbar\xi$ modulo a smooth
  function, where $\chi\in\CI(\R)$, $\chi(\xi)=1$ for $\xi\ge1$ and
  $\chi(\xi)=0$ for $\xi\le-1$. Thus the amplitude vanishes identically in the
  negative direction.
\end{theorem}
\begin{proof}
  The kernel of $S$ with respect to $d\theta'$ is the distribution
  \begin{equation*}
    K_S(\theta,\theta')=\frac1{2\pi}\sum_{k\ge0}e^{ik(\theta-\theta')}
    =\lim_{r\uparrow1}\frac1{2\pi}\frac1{1-re^{i\phi}},\quad
    \phi=\theta-\theta'.
  \end{equation*}
  The limit is taken in distributions, since the Fourier coefficients
  $r^k\to1$ with uniform polynomial bounds. For $\phi\notin2\pi\Z$ the limit
  is the smooth function $(2\pi)^{-1}(1-e^{i\phi})^{-1}$, locally uniformly
  with all derivatives. So $K_S$ is smooth off the diagonal. Near $\phi=0$,
  write $r=e^{-\epsilon}$, so $re^{i\phi}=e^{i(\phi+i\epsilon)}$, and use the
  Laurent expansion
  \begin{equation*}
    \frac1{1-e^{iz}}=\frac iz+h(z),\quad h\text{ holomorphic in }\abs z<2\pi
    \quad(h(0)=\tfrac12).
  \end{equation*}
  Hence for $\abs\phi<\pi$
  \begin{equation*}
    K_S=\frac1{2\pi}\Bigl(\frac i{\phi+i0}+h(\phi)\Bigr),\qquad
    \frac{i}{\phi+i0}=\lim_{\epsilon\downarrow0}\frac i{\phi+i\epsilon}
    =\lim_{\epsilon\downarrow0}\int_0^\infty e^{i\phi\xi-\epsilon\xi}\,d\xi.
  \end{equation*}
  So $(2\pi)^{-1}i/(\phi+i0)=\int e^{i\phi\xi}H(\xi)\,\dbar\xi$ with $H$ the
  Heaviside function. Since $H-\chi$ has compact support,
  $\int e^{i\phi\xi}(H-\chi)\,\dbar\xi$ is entire in $\phi$. Thus near the
  diagonal $K_S$ is the kernel of $\Op(\chi)$ modulo a smooth function. As
  $\chi$ is a classical symbol of order $0$ (it is homogeneous of degree $0$
  for $\abs\xi\ge1$), $S\in\Psido0(\Sphere^1)$ with principal symbol $\chi$
  restricted to $\{\xi=\pm1\}$.
\end{proof}

\subsection{The Toeplitz algebra}
On $\C^N$-valued functions $S$ acts componentwise. Set
\begin{equation*}
  \PsiT{m}(\Sphere^1;\C^N)=\{A\in\Psido m(\Sphere^1;\C^N);\ A=SAS\},\quad
  m\in\{0,-\infty\},
\end{equation*}
regarded as operators on $\Hardy(\Sphere^1;\C^N)$. This is a subalgebra with
identity $S$: if $A_i=SA_iS$ then
$S(A_1A_2)S=(SA_1S)(SA_2S)=A_1A_2$. Its elements map $\Hardym$ to zero and
$\CI$ into $\Hardy$, and $\PsiT{-\infty}=S\Psido{-\infty}S$. Since
$\sigma_0(A)=\sigma_0(S)\sigma_0(A)\sigma_0(S)$, the principal symbol of
$A\in\PsiT0$ vanishes on $S^*_-\Sphere^1$. We call $A$ \emph{elliptic} (as a
Toeplitz operator) if $\sigma_0(A)|_{S^*_+\Sphere^1}$, a smooth function on
$\Sphere^1$ with values in $M(N,\C)$, is invertible. The basic examples are the
\emph{Toeplitz operators} $T_a=SaS$ with multiplier
$a\in\CI(\Sphere^1;M(N,\C))$. They are elliptic exactly when $a$ takes values
in $\GL(N,\C)$.

\begin{proposition}\label{index:prop:toeplitz-fredholm}
  If $A\in\PsiT0(\Sphere^1;\C^N)$ is elliptic, then $A+(\Id-S)\in\Psido0
  (\Sphere^1;\C^N)$ is elliptic, and $A$ is Fredholm on
  $\Hardy(\Sphere^1;\C^N)$ with
  \begin{equation*}
    \ind_{\mathcal T}(A)=\dim\Nul(A|_{\Hardy})-\dim\bigl(\Hardy/A\Hardy\bigr)
    =\ind\bigl(A+\Id-S\bigr).
  \end{equation*}
  There is $B\in\PsiT0$ with $AB-S$ and $BA-S$ in $\PsiT{-\infty}$. If $A$ is
  invertible on $\Hardy$, its inverse lies in $\PsiT0$. The Toeplitz index
  is constant along smooth families of elliptic elements of $\PsiT0$.
\end{proposition}
\begin{proof}
  The symbol of $P=A+\Id-S$ is $\sigma_0(A)$ on $S^*_+$ and $\Id$ on $S^*_-$,
  so $P$ is elliptic. With respect to $\CI=\Hardy\oplus\Hardym$, $P$ is
  diagonal, acting as $A$ on $\Hardy$ and as $\Id$ on $\Hardym$. So the null
  space and cokernel of $P$ on $\CI$ are those of $A$ on $\Hardy$. If $Q$ is a
  parametrix (or the inverse) of $P$, then $B=SQS$ satisfies
  $BA=SQSAS=SQPS=S-S(\Id-QP)S$. The other identity is similar. The last
  statement follows from \cref{index:thm:homotopy} applied to the smooth
  family $A_t+\Id-S$.
\end{proof}

\subsection{The index on the circle}
The \emph{shifts} on $\Hardy(\Sphere^1)$ are
\begin{equation*}
  U=Se^{i\theta}S:\sum_{k\ge0}c_ke^{ik\theta}\mapsto\sum_{k\ge0}c_ke^{i(k+1)\theta},
  \quad
  L=Se^{-i\theta}S:\sum_{k\ge0}c_ke^{ik\theta}\mapsto\sum_{k\ge1}c_ke^{i(k-1)\theta}.
\end{equation*}
Both are elliptic elements of $\PsiT0(\Sphere^1)$, and
\begin{equation}\label{index:eq:shifts}
  LU=\Id,\quad UL=\Id-\pi_0,
\end{equation}
where $\Id$ means $S$, the identity on $\Hardy$, and $\pi_0$ is the orthogonal
projection onto the constants. Thus $U$ is injective with range of
codimension one (complemented by the constants), and $L$ is surjective with
null space the constants:
\begin{equation*}
  \ind_{\mathcal T}(U)=-1,\qquad\ind_{\mathcal T}(L)=1.
\end{equation*}
More generally $U^k=T_{e^{ik\theta}}$ for $k\ge0$ has index $-k$ and
$L^k=T_{e^{-ik\theta}}$ has index $k$. \Cref{index:fig:shifts} shows the
mechanism.

\begin{figure}[ht]
\centering
\begin{tikzpicture}[>=Stealth]
\begin{scope}
\foreach \k in {0,...,4}{\fill (\k*1.1,0) circle (2.2pt);
  \node[below,font=\footnotesize] at (\k*1.1,-0.1){$\k$};}
\node at (5.3,0) {$\cdots$};
\foreach \k in {0,...,3}{\draw[->,shorten >=3pt,shorten <=3pt]
  (\k*1.1,0) to[bend left=45] (\k*1.1+1.1,0);}
\draw[->,shorten >=3pt,shorten <=3pt] (4.4,0) to[bend left=45] (5.2,0.1);
\draw (0,0) circle (5pt);
\node[font=\footnotesize,align=center] at (0,-0.85) {not hit:\\ $\Hardy/U\Hardy=\C$};
\node[font=\small] at (2.2,1.0) {$U$: $k\mapsto k+1$, \ $\ind_{\mathcal T}U=-1$};
\end{scope}
\begin{scope}[shift={(7.2,0)}]
\foreach \k in {0,...,4}{\fill (\k*1.1,0) circle (2.2pt);
  \node[below,font=\footnotesize] at (\k*1.1,-0.1){$\k$};}
\node at (5.3,0) {$\cdots$};
\foreach \k in {1,...,4}{\draw[->,shorten >=3pt,shorten <=3pt]
  (\k*1.1,0) to[bend right=45] (\k*1.1-1.1,0);}
\draw[->,shorten <=3pt] (0,0) -- (0,0.6) node[above,font=\footnotesize]{$0$};
\node[font=\footnotesize,align=center] at (0,-0.85) {$\Nul L=\C$};
\node[font=\small] at (2.6,1.0) {$L$: $k\mapsto k-1$, \ $\ind_{\mathcal T}L=1$};
\end{scope}
\end{tikzpicture}
\caption{The shifts on the Hardy space, acting on the Fourier modes
$e^{ik\theta}$, $k\ge0$. The forward shift $U$ is injective but misses the
constants. The backward shift $L$ is surjective but kills the constants. On the
full space $\CI(\Sphere^1)$, where $k$ runs over $\Z$, both are invertible.}
\label{index:fig:shifts}
\end{figure}

\begin{question}
  The shift $U$ is an isometry and its symbol $e^{i\theta}$ never vanishes. How
  can it fail to be invertible?
\end{question}
\begin{answer}
  An injective linear map of a finite dimensional space to itself is surjective,
  but this is false in infinite dimensions (``Hilbert's hotel''). Ellipticity
  only says that $U$ is invertible modulo smoothing operators, and the index
  measures the remaining defect. Multiplication by $e^{i\theta}$ on all of
  $\CI(\Sphere^1)$ is invertible; it is the restriction to half of the
  frequencies, i.e.\ to half of $S^*\Sphere^1$, that creates the index.
\end{answer}

\begin{theorem}\label{index:thm:toeplitz-index}
  If $a\in\CI(\Sphere^1;\GL(N,\C))$ then the Toeplitz operator $T_a=SaS$ on
  $\Hardy(\Sphere^1;\C^N)$ is Fredholm and
  \begin{equation*}
    \ind_{\mathcal T}(T_a)=-\wind(\det a)
    =-\frac1{2\pi i}\int_{\Sphere^1}\tr\bigl(a^{-1}da\bigr).
  \end{equation*}
\end{theorem}
\begin{proof}
  The second equality is the computation in the proof of
  \cref{index:prop:integrality}. Let $k=\wind(\det a)$. We claim that $a$ is
  smoothly homotopic in $\CI(\Sphere^1;\GL(N,\C))$ to
  $a_k=\diag(e^{ik\theta},1,\dots,1)$. Write $a=\diag(\det a,1,\dots,1)\,b$
  with $b$ a loop in $\mathrm{SL}(N,\C)$. The scalar loop $\det a$ is
  homotopic in $\C^*$ to $e^{ik\theta}$: write $\det
  a=e^{ik\theta}e^{\phi(\theta)}$ with $\phi$ a smooth function on the
  circle, and deform $\phi$ to $0$. Since $\mathrm{SL}(N,\C)$ is connected and
  simply connected (a standard fact; it retracts onto $\mathrm{SU}(N)$), $b$ is
  smoothly homotopic to the constant loop $\Id$. The family
  $T_{a_t}$ along the homotopy is a smooth family of elliptic Toeplitz
  operators, so by \cref{index:prop:toeplitz-fredholm}
  $\ind_{\mathcal T}(T_a)=\ind_{\mathcal T}(T_{a_k})$. Now
  $T_{a_k}=U^k\oplus\Id\oplus\dots\oplus\Id$ if $k\ge0$ and
  $L^{-k}\oplus\Id\oplus\dots\oplus\Id$ if $k<0$. Its index is $-k$.
\end{proof}

\begin{discussion}
  The proof follows a pattern that recurs in the rest of the book. First,
  homotopy invariance reduces the problem to a symbol in a \emph{normal form},
  here $\diag(e^{ik\theta},1,\dots,1)$. Second, the index of an explicit operator
  with that symbol is computed by hand, here by counting Fourier modes. In
  \cref{ch:bott} the same two steps are carried out for families, where the
  simple homotopy above is replaced by Atiyah's linearization of loops.
\end{discussion}

\begin{remark}\label{index:rem:toeplitz-K}
  Via finite rank embeddings $\GL(N,\C)\hookrightarrow\Ginf$,
  \cref{index:thm:toeplitz-index} says that the Toeplitz index is the
  homomorphism $-U_1/2\pi i:K^{-1}(\Sphere^1)\to\Z$. It vanishes on the
  classes of constant maps, so by \cref{index:lem:K1-split} it factors through
  $K^{-2}(\mathrm{pt})$. By \cref{index:prop:K2-loop} it is there given by
  $-u_0/2\pi i$ on $\pi_0(\Gloop)$. That it is an isomorphism is the case
  $k=1$ of \eqref{index:eq:homotopy-groups}. This is the prototype of the
  index formula.
\end{remark}

\section{Smoothing perturbations of elliptic families}\label{index:sec:families}

An elliptic operator of index zero becomes invertible after adding a suitable
finite rank smoothing operator. For a family $A(x)$ we would like to do this
smoothly in $x$. Index zero at each point is necessary but not sufficient: the
cokernels can twist into a non-trivial vector bundle as $x$ varies. This section
shows that this twisting, measured by the families index, is the only
obstruction.

Throughout this section $X$ is a compact manifold (the parameter space,
possibly with boundary). We use the standard $L^2$ structure on
$\Sphere^1$ and the projections onto finite Fourier series
\begin{equation*}
  \pi_r\Bigl(\sum_kc_ke^{ik\theta}\Bigr)=\sum_{\abs k\le r}c_ke^{ik\theta},
  \quad r\in\N_0,
\end{equation*}
acting componentwise on $\C^N$-valued functions. They are self-adjoint finite
rank smoothing operators of rank $(2r+1)N$.

\begin{lemma}\label{index:lem:pi-r}
  If $R\in\CI(X;\Psido{-\infty}(\Sphere^1;\C^N))$ then $R\pi_r\to R$ and
  $\pi_rR\to R$ in $\CI(X;\Psido{-\infty})$, uniformly on $X$, as
  $r\to\infty$.
\end{lemma}
\begin{proof}
  The kernel of $R_x\pi_r$ is
  \begin{equation*}
    \int R_x(\theta,\theta'')\pi_r(\theta'',\theta')\,d\theta''=
    \sum_{\abs k\le r}b_k(x,\theta)e^{-ik\theta'},\quad
    b_k(x,\theta)=\frac1{2\pi}\int R_x(\theta,\theta'')e^{ik\theta''}\,d\theta'',
  \end{equation*}
  i.e.\ the partial sums of the Fourier series of $R_x(\theta,\theta')$ in the
  second variable. Integration by parts in $\theta''$ shows that every
  $(x,\theta)$-derivative of $b_k$ is $O(\ang k^{-q})$ for all $q$, uniformly
  in $(x,\theta)$. So the series converges with all derivatives. The case of
  $\pi_rR$ follows by taking adjoints.
\end{proof}

\begin{lemma}\label{index:lem:stabilized}
  Let $A\in\CI(X;\Psido0(\Sphere^1;\C^N))$ be a smooth family of elliptic
  operators and $A'_r(x)=A(x)(\Id-\pi_r)$. There is $r_0$ such that for all
  $r\ge r_0$ and $x\in X$:
  \begin{enumerate}[label=(\alph*)]
  \item $\Nul(A'_r(x))=\Ran(\pi_r)$, of dimension $(2r+1)N$;
  \item the range of $A'_r(x)$ is closed, with complement
    \begin{equation*}
      F_r(x)=\Nul\bigl((\Id-\pi_r)A(x)^*\bigr)=\Ran(A'_r(x))^\perp
    \end{equation*}
    of dimension $(2r+1)N-\ind A(x)$;
  \item the $F_r(x)$ form a smooth vector bundle $F_r\to X$.
  \end{enumerate}
\end{lemma}
\begin{proof}
  Let $Q$ be a smooth family of parametrices,
  $Q(x)A(x)=\Id-R(x)$ with $R\in\CI(X;\Psido{-\infty})$
  (\cref{index:lem:family-parametrix}). Then
  \begin{equation*}
    Q(x)A'_r(x)=(\Id-R(x))(\Id-\pi_r)=\bigl(\Id-R(x)(\Id-\pi_r)\bigr)(\Id-\pi_r).
  \end{equation*}
  By \cref{index:lem:pi-r}, $R(x)(\Id-\pi_r)\to0$ uniformly in
  $\CI(X;\Psido{-\infty})$, hence in $L^2$ norm. So for $r\ge r_0$ the
  operator $\Id-R(x)(\Id-\pi_r)$ is invertible, with inverse $\Id-S_r(x)$
  depending smoothly on $x$ in $\Psido{-\infty}$ by (vi). With
  $Q'_r=(\Id-S_r)Q$ this gives
  \begin{equation}\label{index:eq:Qprime}
    Q'_r(x)A'_r(x)=\Id-\pi_r.
  \end{equation}
  So $\Nul(A'_r)\subset\Ran(\pi_r)$, and the reverse inclusion is clear, which
  proves (a). $A'_r(x)$ is elliptic and differs from $A(x)$ by a smoothing
  operator, so it has the same index. By
  \cref{index:prop:generalized-inverse} its range has the orthogonal
  complement $\Nul((A'_r)^*)=F_r(x)$, of dimension $(2r+1)N-\ind A(x)$. This
  proves (b).

  For (c), the dimension is locally constant by \cref{index:thm:homotopy}. Let
  $T(x)=A'_r(x)^*A'_r(x)+\pi_r$, a smooth family of elliptic self-adjoint
  operators of order $0$. It is injective, since $T(x)u=0$ implies
  $\norm{A'_ru}^2+\norm{\pi_ru}^2=0$, so $u\in\Ran(\pi_r)\cap\Nul(\pi_r)=\{0\}$.
  It is therefore invertible by \eqref{index:eq:range}, and
  $G(x)=T(x)^{-1}$ is a smooth family by
  \cref{index:lem:family-parametrix}(b). Set $P(x)=A'_rGA_r'^*$. Since
  $A'_r\pi_r=0$ we have $T\pi_r=\pi_r$, hence $G\pi_r=\pi_r$ and
  $GA_r'^*A'_r=G(T-\pi_r)=\Id-\pi_r$. Therefore
  $P^2=A'_r(\Id-\pi_r)GA_r'^*=P$, $P^*=P$ and $PA'_r=A'_r(\Id-\pi_r)=A'_r$.
  Since also $\Ran P\subset\Ran A'_r$, $P(x)$ is the orthogonal projection
  onto $\Ran(A'_r(x))$. Thus $\Id-P(x)$ is the orthogonal projection onto $F_r(x)$. It
  is a smooth family in $\Psido0$ whose values are smoothing (finite rank with
  range in $\CI$). Its full symbol is therefore residual for each $x$, which
  for a smooth family of classical symbols means the family lies in
  $\CI(X;\Psido{-\infty})$. If $v_1,\dots,v_d$ is a basis of $F_r(x_0)$, then
  $(\Id-P(x))v_i$ are smooth sections of $F_r$ which remain a basis near
  $x_0$. So $F_r$ is a smooth vector bundle.
\end{proof}

The formal difference ``$\Ran(\pi_r)\ominus F_r$'' of the trivial bundle of
null spaces and the bundle $F_r$ complementary to the range is the
\emph{families index} of $A$, an element of $K^0(X)$. It is discussed in
\cref{ch:families}.

\begin{discussion}
  Why cut off the first $2r+1$ Fourier modes? The null spaces $\Nul A(x)$ may
  jump in dimension as $x$ varies, so they do not form a bundle. Composing with
  $\Id-\pi_r$ forces the null space to be the \emph{fixed} space $\Ran(\pi_r)$
  for every $x$ (\cref{index:lem:stabilized}(a)). All the variation is then
  pushed into the complement $F_r(x)$ of the range, whose dimension is constant
  because the index is. The trivial bundle $\Ran(\pi_r)$ and the bundle $F_r$
  together carry exactly the information that a single number, the index,
  carries for one operator.
\end{discussion}

\begin{theorem}\label{index:thm:perturb-family}
  Let $A\in\CI(X;\Psido0(\Sphere^1;\C^N))$ be a smooth family of elliptic
  operators. There is a smooth family $B\in\CI(X;\Psido{-\infty}(\Sphere^1;\C^N))$
  such that $A(x)+B(x)$ is invertible for every $x\in X$ if and only if, for
  some (equivalently every) $r\ge r_0$, the bundle $F_r$ of
  \cref{index:lem:stabilized} is isomorphic to the trivial bundle
  $X\times\C^{(2r+1)N}$.
\end{theorem}
\begin{proof}
  Suppose $F_r$ is trivial for some $r\ge r_0$. Since $\Ran(\pi_r)$ has
  dimension $(2r+1)N$, there is a smooth bundle isomorphism
  $\Phi(x):\Ran(\pi_r)\to F_r(x)$. Then $B_r(x)=\Phi(x)\pi_r$ is a smooth
  family of finite rank smoothing operators: its kernel is
  $\sum_j\Phi(x)e_j\otimes e_j^*$ for an orthonormal basis $e_j$ of
  $\Ran(\pi_r)$. The operator $A'_r(x)+B_r(x)$ maps $\Ran(\Id-\pi_r)$
  injectively onto $\Ran(A'_r(x))$ and $\Ran(\pi_r)$ isomorphically onto
  $F_r(x)$. As $L^2=\Ran(A'_r)\oplus F_r$, it is invertible. Thus
  $B=B_r-A\pi_r$ has the required property, since $A+B=A'_r+B_r$.

  Conversely, suppose $A+B$ is invertible on $X$. By \cref{index:lem:pi-r},
  $B\pi_r\to B$ in $L^2$ norm, uniformly on $X$. Since the inverses
  $(A(x)+B(x))^{-1}$ are bounded uniformly on the compact set $X$,
  $A+B\pi_r$ is invertible for all large $r$. Write
  \begin{equation*}
    A+B\pi_r=A'_r+B_r,\quad B_r=(A+B)\pi_r.
  \end{equation*}
  Invertibility of $A'_r+B_r$ means that
  $L^2=A'_r(\Ran(\Id-\pi_r))\oplus B_r(\Ran(\pi_r))$. So
  $V(x)=B_r(x)\Ran(\pi_r)$ is a smooth bundle complementary to
  $\Ran(A'_r(x))$, trivialized by $B_r$ (which is injective on
  $\Ran(\pi_r)$). The orthogonal projection $\Id-P(x)$ onto $F_r(x)$ restricts
  to an isomorphism $V(x)\to F_r(x)$, since its null space is $\Ran(A'_r(x))$.
  Hence $F_r$ is trivial for all large $r$. Together with the first part,
  triviality for one $r\ge r_0$ implies triviality for all large $r$.
\end{proof}

In particular a necessary condition is that $\ind A(x)=0$, and when $X$ is a
point the theorem reduces to the statement of \cref{ch:manifolds} that an
elliptic operator has an invertible smoothing perturbation if and only if its
index vanishes. The proofs extend to any compact manifold in place of the
circle (and to fibrations by compact manifolds) once a sequence of finite rank
smoothing projections with the property of \cref{index:lem:pi-r} is
available. They also pass unchanged to the Toeplitz algebra if $\pi_r$ is
replaced by the projection $\pi^+_r$ onto $\operatorname{span}\{e^{ik\theta};\
0\le k\le r\}$. This projection commutes with $S$, lies in $\PsiT{-\infty}$,
and has rank $(r+1)N$.

\begin{corollary}\label{index:cor:toeplitz-perturb}
  Let $A\in\CI(X;\PsiT0(\Sphere^1;\C^N))$ be a smooth family of elliptic
  Toeplitz operators. There is $B\in\CI(X;\PsiT{-\infty}(\Sphere^1;\C^N))$
  with $A(x)+B(x)$ invertible on $\Hardy$ for all $x$ if and only if, for
  large $r$, the bundle $F^+_r(x)=\Nul((S-\pi^+_r)A(x)^*)\cap\Hardy$ is
  trivial of rank $(r+1)N$. The inverses then lie in $\PsiT0$.
\end{corollary}

\begin{corollary}\label{index:cor:lift-homotopy}
  Let $A\in\CI([0,1]\times X;\PsiT0(\Sphere^1;\C^N))$ (or
  $\CI([0,1]\times X;\Psido0(\Sphere^1;\C^N))$) be elliptic at each point, with
  $A_0(x)$ invertible for all $x\in X$. Then there is a smooth family
  $B\in\CI([0,1]\times X;\PsiT{-\infty})$ (respectively $\Psido{-\infty}$) with
  $B_0=0$ and $A_t(x)+B_t(x)$ invertible for all $(t,x)$.
\end{corollary}
\begin{proof}
  We give the proof in the $\Psido0$ case; the Toeplitz case is identical. Let
  $r$ be large. Applying the converse half of
  \cref{index:thm:perturb-family} on $X$ with $B=0$ shows that the
  restriction of $F_r$ to $\{0\}\times X$ is trivial. Since $[0,1]\times X$
  retracts onto $\{0\}\times X$, the bundle $F_r$ over $[0,1]\times X$ is
  isomorphic to the pull-back of its restriction to $\{0\}\times X$ (homotopy
  invariance of vector bundles \cite{Atiyah-K}), so it is trivial. The first
  half of the theorem gives a smooth family $B'$ of smoothing operators with
  $C_t=A_t+B'_t$ invertible everywhere. To make the perturbation vanish at
  $t=0$, set
  \begin{equation*}
    D_t=C_tC_0^{-1}A_0=A_t+B_t,\quad
    B_t=-A_t(A_0+B'_0)^{-1}B'_0+B'_tC_0^{-1}A_0.
  \end{equation*}
  Here we used $C_0^{-1}A_0=\Id-C_0^{-1}B'_0$. $D_t$ is invertible as a product
  of invertibles, $D_0=A_0$, and $B_t$ is a smooth family of smoothing
  operators, since $C_0^{-1}\in\Psido0$ by
  \cref{index:lem:family-parametrix}(b). It vanishes at $t=0$ because
  $D_0=A_0$.
\end{proof}

This allows homotopies of elliptic families to be lifted to homotopies of
invertible families. It is used in \cref{ch:bott}.

\section{Deforming matrices away in the Toeplitz group}\label{index:sec:deform}

In \cref{ch:bott} we show that the group of invertible Toeplitz operators is
(weakly) contractible. A first obstacle is the copy of $\GL(N,\C)$ inside it:
the constant matrices, which are not contractible in $\GL(N,\C)$ itself. The
point of this section is that, after doubling $\C^N$ to $\C^{2N}$, the shifts
$U$ and $L$ provide enough room to contract them.

Let
\begin{equation*}
  \tilde G^0_{\mathcal T}(\Sphere^1;\C^N)=\{A\in\PsiT0(\Sphere^1;\C^N);\
  A\text{ is elliptic and invertible on }\Hardy,\ A^{-1}\in\PsiT0\}
\end{equation*}
(the tilde distinguishes it from a smaller group used in \cref{ch:bott}). By
\cref{index:prop:toeplitz-fredholm} the last condition is automatic. The
smoothing elements $G^{-\infty}_{\mathcal T}=\tilde G^0_{\mathcal
T}\cap(S+\PsiT{-\infty})$ form a subgroup containing a copy of $\GL(N,\C)$:
\begin{equation}\label{index:eq:GL-embed}
  \GL(N,\C)\ni g\longmapsto\Id-\pi_0+\pi_0g\pi_0\in
  G^{-\infty}_{\mathcal T}(\Sphere^1;\C^N),
\end{equation}
which acts as $g$ on the constants and as the identity on higher modes; here
$\Id=S$. We also use the standard stabilization
$\GL(N,\C)\hookrightarrow\GL(2N,\C)$,
$g\mapsto\left(\begin{smallmatrix}g&0\\0&\Id_N\end{smallmatrix}\right)$.

\begin{proposition}\label{index:prop:deform-GL}
  Embed $\GL(N,\C)$ in $\tilde G^0_{\mathcal T}(\Sphere^1;\C^{2N})$ by
  combining the stabilization with \eqref{index:eq:GL-embed} (for $2N$ in
  place of $N$). Then there is a continuous map
  $M:[0,1]\times\GL(N,\C)\to\tilde G^0_{\mathcal T}(\Sphere^1;\C^{2N})$,
  smooth in all variables, such that $M(0,g)$ is the image of $g$ and
  $M(1,g)=\Id$ for all $g$.
\end{proposition}
\begin{proof}
  Write $\C^{2N}=\C^N\oplus\C^N$ and operators as $2\times2$ block matrices of
  $N\times N$ Toeplitz operators. A constant matrix $g$ commutes with $S$, so
  it defines the Toeplitz operator $gS$, which we write as $g$. The image of
  $g$ is
  \begin{equation*}
    M_0=\begin{pmatrix}\Id-\pi_0+\pi_0g\pi_0&0\\0&\Id\end{pmatrix}.
  \end{equation*}
  The homotopy has three stages, each depending smoothly on $g$:
  \begin{equation*}
    \begin{tikzcd}[column sep=2.6em]
      M_0 \arrow[r,"\text{Stage 1}"] &
      M_{\pi/2} \arrow[r,"\text{Stage 2}"] &
      \begin{pmatrix}g&0\\0&g^{-1}\end{pmatrix} \arrow[r,"\text{Stage 3}"] &
      \Id .
    \end{tikzcd}
  \end{equation*}
  In Stages 1 and 2 the Toeplitz structure is essential: the shifts move the
  matrix $g$ off the constants, so that it can be traded for $g^{-1}$ in the
  second block. Stage 3 is the finite dimensional fact that
  $\diag(g,g^{-1})$ is connected to $\Id$ uniformly in $g$.

  \emph{Stage 1.} For $\tau\in[0,\pi/2]$ let
  \begin{equation*}
    M_\tau=\begin{pmatrix}\cos\tau(\Id-\pi_0)+\pi_0g\pi_0&\sin\tau\,gU\\
    -\sin\tau\,g^{-1}L&\cos\tau\,\Id\end{pmatrix}.
  \end{equation*}
  Its symbol on $S^*_+$ is
  $\left(\begin{smallmatrix}\cos\tau&\sin\tau\,ge^{i\theta}\\
  -\sin\tau\,g^{-1}e^{-i\theta}&\cos\tau\end{smallmatrix}\right)$, of
  determinant $1$, so $M_\tau$ is elliptic. To see that it is invertible,
  write $u=\sum_{k\ge0}u_ke^{ik\theta}$, $v=\sum_{k\ge0}v_ke^{ik\theta}$ with
  $u_k,v_k\in\C^N$. Then $M_\tau$ maps $(u_0,0)$ to $(gu_0,0)$. For $k\ge1$ it
  maps the pair of modes $(u_ke^{ik\theta},v_{k-1}e^{i(k-1)\theta})$ to
  $(f_ke^{ik\theta},h_{k-1}e^{i(k-1)\theta})$ with
  \begin{equation*}
    \begin{pmatrix}f_k\\h_{k-1}\end{pmatrix}=
    \begin{pmatrix}\cos\tau&\sin\tau\,g\\-\sin\tau\,g^{-1}&\cos\tau\end{pmatrix}
    \begin{pmatrix}u_k\\v_{k-1}\end{pmatrix},
  \end{equation*}
  and
  \begin{equation*}
    \begin{pmatrix}\cos\tau&\sin\tau\,g\\-\sin\tau\,g^{-1}&\cos\tau\end{pmatrix}
    =\begin{pmatrix}1&0\\0&g^{-1}\end{pmatrix}R_\tau
    \begin{pmatrix}1&0\\0&g\end{pmatrix},
  \end{equation*}
  where
  $R_\tau=\left(\begin{smallmatrix}\cos\tau&\sin\tau\\-\sin\tau&\cos\tau
  \end{smallmatrix}\right)$. These matrices are invertible, with inverses of
  the same form with $\tau$ replaced by $-\tau$. So $M_\tau$ is invertible, and
  its inverse is the Toeplitz operator obtained by replacing $\tau$ by $-\tau$
  and $g$ by $g^{-1}$ in the $(1,1)$ entry. Thus $M_\tau\in\tilde
  G^0_{\mathcal T}$. At $\tau=\pi/2$,
  \begin{equation*}
    M_{\pi/2}=\begin{pmatrix}\pi_0g\pi_0&gU\\-g^{-1}L&0\end{pmatrix}.
  \end{equation*}

  \emph{Stage 2.} For $\tau$ decreasing from $\pi/2$ to $0$ let
  \begin{equation*}
    M'_\tau=\begin{pmatrix}\cos\tau\,g(\Id-\pi_0)+\pi_0g\pi_0&\sin\tau\,gU\\
    -\sin\tau\,g^{-1}L&\cos\tau\,g^{-1}\end{pmatrix},
  \end{equation*}
  which equals $M_{\pi/2}$ at $\tau=\pi/2$. Its symbol is
  $\diag(g,g^{-1})\left(\begin{smallmatrix}\cos\tau&\sin\tau\,e^{i\theta}\\
  -\sin\tau\,e^{-i\theta}&\cos\tau\end{smallmatrix}\right)$, which is
  invertible. On modes it acts as $g$ on $u_0$ and by the matrix
  $\left(\begin{smallmatrix}\cos\tau\,g&\sin\tau\,g\\-\sin\tau\,g^{-1}&
  \cos\tau\,g^{-1}\end{smallmatrix}\right)=\diag(g,g^{-1})R_\tau$ on the pairs
  $(u_k,v_{k-1})$. So it is again in $\tilde G^0_{\mathcal T}$. At $\tau=0$
  we reach the constant matrix
  \begin{equation*}
    M'_0=\begin{pmatrix}g&0\\0&g^{-1}\end{pmatrix}.
  \end{equation*}

  \emph{Stage 3.} It remains to deform $\diag(g,g^{-1})$ to the identity in
  $\GL(2N,\C)$, uniformly in $g$. First take
  $\left(\begin{smallmatrix}\cos\tau\,g&\sin\tau\\-\sin\tau&\cos\tau\,g^{-1}
  \end{smallmatrix}\right)$, $\tau\in[0,\pi/2]$. For $\cos\tau\ne0$ the Schur
  complement of the invertible $(1,1)$ entry is
  $\cos\tau\,g^{-1}+\sin^2\tau(\cos\tau)^{-1}g^{-1}=(\cos\tau)^{-1}g^{-1}$, so
  the matrix is invertible. At $\tau=\pi/2$ it is
  $\left(\begin{smallmatrix}0&1\\-1&0\end{smallmatrix}\right)$. Then follow
  the rotations $R_\tau$ from $\tau=\pi/2$ back to $\tau=0$. Constant
  invertible matrices are elements of $\tilde G^0_{\mathcal T}$. After
  smoothing the corners by reparametrization, the concatenation of the three
  stages is the required homotopy.
\end{proof}

\begin{remark}\label{index:rem:deform-meaning}
  \Cref{index:prop:deform-GL} shows that the ``$G^{-\infty}_{\mathcal T}$
  part'' of the Toeplitz group can be (weakly) deformed away after
  stabilization. In \cref{ch:bott} this is combined with Atiyah's
  linearization of loops in $\GL(N,\C)$, which shows how invertible elliptic
  Toeplitz operators can be deformed into the smoothing subgroup. Together
  these give the weak contractibility of the middle group in
  \eqref{index:eq:classifying}.
\end{remark}

\section{Exercises}

\begin{exercise}\label{index:ex:multiplicative}
  Let $A_1\in\Psido{m_1}(X;F,G)$ and $A_2\in\Psido{m_2}(X;E,F)$ be elliptic.
  Using \cref{index:thm:trace-formula} with the parametrix $B_2B_1$ of
  $A_1A_2$, show that $\ind(A_1A_2)=\ind(A_1)+\ind(A_2)$. (Write
  $\Id-B_2B_1A_1A_2=(\Id-B_2A_2)+B_2R_1A_2$ with $R_1=\Id-B_1A_1$, and use
  \cref{index:lem:trace-commute}.)
\end{exercise}

\begin{exercise}\label{index:ex:adjoint}
  Show that $\ind(A^*)=-\ind(A)$ for elliptic $A$, and that every elliptic
  self-adjoint $A\in\Psido m(X;E)$ has index zero.
\end{exercise}

\begin{exercise}\label{index:ex:same-symbol-family}
  Let $A\in\CI(X';\Psido m(X;E,F))$ be a smooth family of elliptic
  operators parametrized by a connected compact manifold $X'$. Show directly
  from \cref{index:thm:trace-formula} and
  \cref{index:lem:family-parametrix} that $x\mapsto\ind A(x)$ is constant.
\end{exercise}

\begin{exercise}\label{index:ex:det-properties}
  Show that for $a\in\Psido{-\infty}(Y;E)$ the function
  $z\mapsto\det(\Id+za)$ is entire, and that its zeros are the points
  $z=-1/\lambda$ with $\lambda\ne0$ an eigenvalue of $a$. Show further that the
  order of the zero at $-1/\lambda$ is the algebraic multiplicity of
  $\lambda$. (Use the local formula in the proof of
  \cref{index:thm:det-extension} with $W$ containing the generalized
  eigenspace.)
\end{exercise}

\begin{exercise}\label{index:ex:U1-circle}
  Let $f:\Sphere^1\to\Ginf(\Sphere^1)$,
  $f(e^{is})=\Id+(e^{is}-1)\pi_0$, where $\pi_0$ is the projection onto the
  constants. Compute $\int_{\Sphere^1}f^*u_1$ and $\det f(e^{is})$. Deduce
  that $U_1:K^{-1}(\Sphere^1)\to H^1(\Sphere^1;\C)\cong\C$ has image
  $2\pi i\Z$.
\end{exercise}

\begin{exercise}\label{index:ex:even-vanish}
  Show that the forms $\Tr((A^{-1}dA)^{2k})$ on $\Ginf$ vanish identically,
  directly from the antisymmetrized formula \eqref{index:eq:chern-forms} with
  $2k$ entries. Show also that $u_{2k-1}$ is invariant under right
  multiplication, using $R_B^*(A^{-1}dA)=B^{-1}(A^{-1}dA)B$.
\end{exercise}

\begin{exercise}\label{index:ex:commutator}
  Show that $[S,f]\in\Psido{-\infty}(\Sphere^1)$ for every
  $f\in\CI(\Sphere^1)$. Use the symbol $\chi$ of \cref{index:thm:szego} and
  the composition formula, noting that all $\xi$-derivatives of $\chi$ have
  compact support. Deduce that $T_aT_b-T_{ab}\in\PsiT{-\infty}$ for
  $a,b\in\CI(\Sphere^1;M(N,\C))$. Also show that, if $a$ and $b$ are
  trigonometric polynomials, $T_aT_b-T_{ab}$ has finite rank.
\end{exercise}

\begin{exercise}\label{index:ex:valuations}
  A compact manifold $Z$ is recovered from the ring $\CI(Z)$ as the set of
  its multiplicative linear functionals (``valuations'') $p$, with
  $p(fg)=p(f)p(g)$ and $p(1)=1$. Show that the continuous multiplicative
  linear functionals on the ring $\Hardy(\Sphere^1)$ are exactly
  $f\mapsto u_f(z)$, $z\in\overline{\mathbb D}$, where $u_f$ is the
  holomorphic extension of \cref{index:prop:hardy}. So the ``space'' of the
  Hardy ring is the closed disc, not half of the circle.
\end{exercise}

\begin{exercise}\label{index:ex:toeplitz-product}
  Using \cref{index:thm:toeplitz-index}, show that
  $\ind_{\mathcal T}(T_aT_b)=\ind_{\mathcal T}(T_a)+\ind_{\mathcal T}(T_b)$
  for $a,b\in\CI(\Sphere^1;\GL(N,\C))$. Check this against
  \cref{index:ex:multiplicative} applied to $T_a+\Id-S$ and $T_b+\Id-S$.
\end{exercise}

\section*{Notes}
The source lectures stated the existence of smooth families of parametrices,
the smoothness of the bundles $F_r$, and the extension of the Fredholm
determinant to all of $\Psido{-\infty}$ without proof. Here they are supplied
in \cref{index:lem:family-parametrix}, \cref{index:lem:stabilized} and
\cref{index:thm:det-extension}. The second stage of the homotopy in
\cref{index:prop:deform-GL} has been corrected so that it starts where the
first stage ends. The parametrization of $\Sphere^1\setminus\{1\}$ in
\cref{index:prop:K2-loop} has been adjusted to omit the base point. The
classical index theorem for Toeplitz operators on the circle
(\cref{index:thm:toeplitz-index}) is due to F.~Noether and to Gohberg and
Krein; it has been added as the simplest instance of the index formula. Loop
groups and their relation with Toeplitz operators are treated at length in
\cite{Pressley-Segal}. The approach via the isotropic algebra on $\R$ is
carried out in the setting of manifolds with boundary in \cite{Melrose-Rochon}.

The Szeg\H{o} projector also exists in higher dimensions. For the boundary
$\Sphere^{2n-1}$ of the unit ball in $\C^n$ (or of a strictly pseudoconvex
domain) the orthogonal projection onto boundary values of holomorphic
functions is a Fourier integral operator of Hermite type. The index of the
Toeplitz operator $T_a=SaS$ with $a\in\CI(\Sphere^{2n-1};\GL(N,\C))$ is a
universal multiple of $\int_{\Sphere^{2n-1}}\tr((a^{-1}da)^{2n-1})$, i.e.\ of
the pairing of the odd Chern form $u_{2n-1}$ with the fundamental class. This
is a theorem of L.~Boutet de Monvel (\emph{On the index of Toeplitz operators
of several complex variables}, Invent.\ Math.\ 50 (1979)). Theorem
\ref{index:thm:toeplitz-index} is the case $n=1$.
